\chapter{分布式计算模型}

分布式计算把计算扩展到多台机器，但同时引入网络通信和故障。共享内存消失后，线程之间的读写变成消息；锁和原子不再直接可用；时间、顺序和失败都变得不可靠。

\section{从共享内存到消息系统}

单机内存访问通常以纳秒计，网络往返通常以微秒到毫秒计。单机 cache miss 已经很贵，跨网络请求更贵。分布式系统要尽量减少同步往返，使用批处理、流水线、数据本地性和异步通信。

网络还可能丢包、重排、延迟抖动或断开。应用层必须处理超时、重试、重复请求和部分失败。一个请求超时，不代表对方没有执行；重试可能造成重复副作用，因此幂等性非常重要。

\section{RPC、延迟和 in-flight 请求}

RPC 把远程调用包装成类似本地函数调用的形式，但它不是本地函数。参数要序列化，消息要进入内核或用户态网络栈，经过网卡和交换机，到达对端后再反序列化和调度执行。一次 RPC 的固定成本远高于一次函数调用，因此分布式系统要避免细粒度频繁 RPC。

批处理、流水线和异步调用是常见优化。批处理减少每条消息固定开销，流水线隐藏往返延迟，异步调用避免线程阻塞。但它们都会增加错误处理复杂度：部分成功、乱序返回、超时重试和取消都要定义清楚。

RPC 还会隐藏背压问题。本地函数调用如果慢，调用者自然等在当前线程；远程调用如果被异步提交，调用者可能继续制造更多请求，直到队列、连接池或对端服务被压垮。因此 RPC 客户端必须有并发上限、超时、重试预算和排队策略。没有这些边界，异步 RPC 只是把阻塞变成不可控积压。

\topic{时间和顺序}

不同机器没有完美同步时钟。日志时间戳可以帮助观察，但不能单独作为一致性依据。分布式系统常用逻辑时钟、版本号、任期号和日志索引表达顺序关系。

顺序要区分因果顺序和物理时间顺序。请求 A 触发请求 B，B 因果上发生在 A 之后；但两个机器的墙钟时间可能显示 B 的时间戳更早。协议通常不依赖墙钟证明因果，而是使用请求 ID、版本号、日志索引、任期或向量时钟等逻辑信息。日志时间戳主要用于观测和排查，不应被误用成强一致证据。

\section{超时、重试和幂等}

幂等操作执行一次和执行多次效果相同。分布式系统中，客户端超时后常常不知道服务器是否执行了请求。若请求不是幂等的，重试可能造成重复扣款、重复创建或重复写入。常见做法是给请求分配唯一 ID，服务器记录已处理 ID，并在重复请求时返回同一结果。

幂等设计要落到存储结构。只在客户端生成请求 ID 不够，服务器还要把请求 ID、执行结果和业务状态放在同一个原子提交边界里。否则服务器可能已经修改业务状态，但还没记录请求 ID 就崩溃，重试后仍然重复执行。工程上常用去重表、条件写、事务日志或状态机日志解决这个问题。

\section{故障、降级和背压}

机器可能宕机，进程可能重启，网络可能分区，磁盘可能慢，消息可能重复。多数工程系统采用 crash-stop 或 crash-recovery 模型，不处理任意恶意错误。明确故障模型，是协议设计的第一步。

\topic{数据本地性}

分布式计算中，移动计算到数据附近通常比移动大量数据更划算。MapReduce、Spark、分布式数据库和对象存储都围绕数据分片、任务调度和网络 shuffle 做权衡。

数据本地性是 NUMA 原则在集群尺度上的版本。单机里远端 NUMA 访问慢，集群里跨机拉取数据更慢；单机里先在线程本地 partial 再合并，集群里先在 map 端 combine 再 shuffle；单机里热点 cache line 会拖慢所有核心，集群里热点 key 会拖慢整个 stage。前面章节的思想在这里不是类比游戏，而是同一原则换了成本单位。

\topic{CAP 的正确理解}

CAP 不是说系统只能在一致性、可用性、分区容忍性里任选两个。网络分区是分布式系统必须面对的现实；在分区发生时，系统需要在继续响应请求和保持强一致之间做选择。正常情况下，系统仍然要同时追求低延迟、高可用和足够一致。工程讨论要具体到故障场景，而不是只背缩写。

\topic{超时、重试和重试风暴}

超时不是失败证明，只是本地等待时间超过阈值。请求可能在网络上，可能在对端队列里，可能已经执行但响应丢失，也可能确实失败。重试可以提高成功率，也可能放大过载。如果一个服务已经慢了，所有客户端同时重试，会制造重试风暴，让系统更慢。

健康的重试策略通常包括指数退避、抖动、最大重试次数、总超时预算、幂等保证和熔断。重试预算尤其重要：一个上游请求如果调用十个下游，每个下游都重试三次，总请求量可能迅速膨胀。分布式系统的性能问题常常来自这些控制路径，而不是用户函数本身。

\begin{keyidea}
分布式系统的性能问题往往是正确性问题的影子。重试、复制、超时、共识和恢复都会改变吞吐和延迟。
\end{keyidea}

\topic{从共享内存到消息}

单机多线程里，线程通过共享内存通信。即使共享内存有 cache line、一致性协议和内存序成本，程序仍然可以用一个地址表达共同状态。分布式系统里，这个前提消失了。机器 A 不能直接读取机器 B 的内存；它只能发送消息，请求 B 执行某个操作或返回某份数据。共享变量变成消息，锁变成协议，原子操作变成带版本号或日志索引的提交。

这个变化非常根本。单机中一个 mutex 临界区通常在纳秒到微秒级，跨网络的一次协调可能是微秒到毫秒级，还会遇到排队、丢包、重试和对端故障。单机中线程崩溃通常意味着进程整体失败，分布式中一个节点失败后其他节点可能继续运行。单机中内存序由语言和硬件定义，分布式中消息顺序、重复和超时都要由协议处理。

因此分布式系统不能把远程调用当成本地函数。RPC 库可以隐藏序列化细节，但不能隐藏失败语义。调用者必须问：请求是否到达，对端是否执行，响应是否丢失，重试是否安全，重复请求如何处理，超时后资源如何回收。若这些问题没有答案，系统在正常路径看起来能跑，在故障路径会变得不可控。

\topic{消息的生命周期}

一条消息从客户端到服务器，至少经过序列化、排队、发送、网络传输、接收、反序列化、调度、执行、生成响应、响应传输和客户端处理。每一步都可能排队，每一步都可能失败。性能报告如果只写“RPC 延迟 3 ms”，还不够；要拆出客户端排队、网络 RTT、服务端队列、执行时间和响应排队。

消息还需要身份。没有 request id，客户端无法把响应和请求可靠对应；没有 attempt id，服务端无法区分第一次请求和重试；没有 trace id，观测系统无法串起跨节点路径；没有 partition id，shuffle 系统无法知道数据应写到哪里。系统工程中，消息头不是形式，它承载正确性和观测性。

\begin{lstlisting}[language=C++,caption={本机模拟用的消息元数据}]
struct MessageHeader {
    std::uint64_t request_id = 0;
    std::uint32_t source_worker = 0;
    std::uint32_t target_partition = 0;
    std::uint32_t attempt = 0;
};

struct ShuffleMessage {
    MessageHeader header;
    std::vector<std::int32_t> keys;
    std::vector<std::int64_t> values;
};
\end{lstlisting}

这段代码不是网络协议实现，但它展示了分布式消息需要的基本边界。request id 用于幂等和去重，source worker 用于观测和调试，target partition 用于路由，attempt 用于重试分析。真实系统还会有校验、版本、压缩格式、认证信息和时间戳。

\topic{延迟、吞吐和 in-flight 请求}

分布式系统里，吞吐往往依赖足够的 in-flight 请求。一个请求的往返延迟可能很高，如果客户端每次只发一个请求，吞吐会被 RTT 限制。流水线允许多个请求同时在网络中飞行，提高吞吐。但 in-flight 太多会导致队列积压、内存占用上升、对端过载和尾延迟恶化。

这和 CPU 的内存级并行度很像：核心可以同时挂起多个 cache miss 来隐藏延迟，但 miss 太多会占满 load buffer；客户端可以同时发多个 RPC 来隐藏 RTT，但请求太多会占满队列和连接。系统要有并发窗口，窗口大小要由延迟、吞吐、错误率和对端能力共同决定。

一个简单模型是 Little's Law：在稳定系统中，平均并发数约等于吞吐乘以平均延迟。若希望每秒处理 10000 个请求，平均延迟 10 ms，则系统中大约有 100 个请求在飞行。这个模型不能替代测量，但能帮助你判断一个并发上限是否荒谬。

\topic{超时不是取消}

客户端超时只表示客户端不再等待，不表示服务器停止执行。服务器可能已经完成业务状态修改，只是响应丢失；也可能还在队列中等待；也可能稍后才执行。若客户端超时后立即重试，而操作不是幂等的，就会重复副作用。若客户端释放请求上下文，但底层完成回调稍后到达，就可能访问已释放对象。

因此超时要配合取消和幂等。取消是尽力而为，不一定成功；幂等保证重复执行的外部效果可控；请求上下文保证完成路径无论成功、失败、超时还是取消，都能安全回收。异步 I/O 章节里的生命周期问题，在分布式 RPC 中更严重，因为远端是否执行不是本地进程能完全控制的。

\topic{幂等提交点}

幂等性不是“重试时先查一下”这么简单。真正的幂等需要把请求 ID、业务状态修改和响应结果放在同一个提交边界里。考虑转账、创建订单或写检查点：如果服务器先修改业务状态，随后记录 request id 前崩溃，重试时看不到旧 request id，就可能重复执行；如果先记录 request id，但业务状态没修改就崩溃，重试时又可能误以为已经成功。

一种常见设计是把请求处理当作状态机事务：检查 request id 是否已处理；若已处理，返回保存的结果；若未处理，执行业务修改，同时写入 request id 和结果；提交后再返回。这个过程可以用数据库事务、共识日志或本地 WAL 实现。关键是提交点必须清楚。

\topic{消息顺序和重复}

网络可能重排消息。客户端发送 A 再发送 B，不代表服务器一定先处理 A。即使 TCP 在单连接内保持字节顺序，应用层多连接、多队列、多线程处理仍可能改变完成顺序。协议如果依赖顺序，就必须显式携带序号、版本或日志索引。

消息还可能重复。客户端重试会重复发送，网络代理可能重放，服务端超时后客户端又发下一次 attempt。服务端必须决定重复消息如何处理：直接丢弃、返回旧结果、合并、覆盖还是报错。没有 request id 和幂等表，重复消息通常无法安全处理。

\topic{故障检测的模糊性}

分布式系统无法可靠地区分“节点死了”和“节点很慢”。心跳超时可能因为节点崩溃、网络拥塞、GC 暂停、磁盘卡顿、CPU 被抢占或调度延迟。把慢节点误判为死节点可能导致重复调度和资源浪费；把死节点长期当成活节点会拖慢恢复。故障检测本质上是基于超时的怀疑，不是数学证明。

工程上通常使用多次心跳、租约、任期、健康检查、慢节点隔离和人工可观测指标。对计算任务而言，可以把慢任务 speculative retry 到其他节点，但要处理重复输出；对存储副本而言，不能轻易让旧 leader 继续写，需要任期和共识保证。故障检测和一致性协议紧密相连。

\topic{部分失败和降级}

单机程序常常整体成功或整体失败，分布式系统经常部分失败。一个查询请求可能需要访问十个分片，其中九个成功、一个超时。系统要定义结果：失败整个请求，返回部分结果，使用缓存旧值，还是降级为近似结果。这个选择取决于业务语义。

搜索结果可以返回部分分片并标记结果不完整；计费系统不能忽略一个失败分片；监控系统可以延迟补齐；推荐系统可以使用缓存。分布式教材不能只讲协议，还要让读者看到业务语义如何决定失败处理。

\topic{背压和重试预算}

分布式背压要跨节点传播。服务端队列满时，应让客户端降低发送速率、拒绝低优先级请求或返回明确错误。若服务端已经过载，客户端仍按固定策略重试，只会放大故障。健康的客户端要有最大并发、总超时、重试次数、指数退避、抖动和熔断。

重试预算可以按上游请求计算。例如一个用户请求最多允许产生 20 次下游调用尝试，超过就停止重试并返回错误。这样可以防止多层服务各自重试导致请求量指数膨胀。预算还应被观测：报告中要记录原始请求数、重试请求数、超时数、最终失败数和重试原因。

\topic{数据本地性和调度}

分布式计算里，调度任务时要考虑数据位置。若输入分片在节点 A，本地执行可以避免网络读取；若节点 A 忙或失败，可以远程读取或复制数据，但要付出网络成本。数据本地性不是绝对规则，而是调度决策的一项成本。

这和 NUMA 类似。单机里，线程靠近内存节点更快；集群里，任务靠近数据分片更快。若本地节点排队很长，远程执行可能更快；若远程网络拥塞，本地等待更好。调度器需要把数据位置、队列长度、节点健康、资源类型和任务优先级一起考虑。

\topic{本机多进程模拟}

本册不要求读者立刻搭建真实集群。可以先用本机多进程或多线程模拟分布式系统：每个 worker 有自己的输入分片和消息队列，driver 负责分配任务，网络用有界队列和随机延迟模拟，故障用随机丢弃、超时和重复消息模拟。这样可以在一台机器上训练协议思维。

模拟不能证明真实网络性能，但能验证语义：请求 ID 是否去重，重试是否幂等，队列是否背压，检查点是否能恢复，重复输出是否被提交协议挡住。先在本机把语义做清楚，再迁移到真实网络环境，风险更低。

\topic{RPC 不是函数调用}

远程调用看起来像函数调用，但语义完全不同。本地函数调用要么进入函数体，要么在调用前就失败；调用栈、对象生命周期和异常传播都在同一进程内。RPC 经过序列化、发送队列、内核网络栈、网卡、交换机、对端内核、服务端队列、业务线程、响应路径。任何一段都可能排队、失败、重复、超时或部分完成。把 RPC 当成本地函数，是分布式系统设计里最常见的错误之一。

RPC 的调用者必须接受一个事实：超时后不知道远端到底发生了什么。远端可能没有收到请求，可能收到但还没执行，可能已经执行并提交，可能响应丢失，可能响应正在返回。这个不确定性决定了幂等、请求 ID、重试预算和提交点不是附加功能，而是 RPC 的基本语义。只要系统允许超时重试，就必须回答重复请求如何处理。

服务端也不能把 RPC 当成本地函数。请求进入服务端后，首先进入接收队列，再由某个 worker 处理。若业务处理很慢，接收队列会增长；若队列无限，内存会增长；若队列有限，客户端会遇到拒绝或超时。服务端处理完后，响应也可能在发送队列中等待。一次 RPC 的延迟不是业务函数耗时，而是排队、执行、序列化和网络共同组成的端到端路径。

\topic{消息边界和协议版本}

分布式系统中，消息必须有明确边界。TCP 是字节流，不保留应用层消息边界；一次 \code{write} 不等于对端一次 \code{read}。应用协议需要 framing，例如固定长度头部加变长 body，或者长度前缀，或者分隔符。没有 framing，服务端无法可靠知道一条消息在哪里结束，下一条消息在哪里开始。

消息还需要版本。系统升级时，客户端和服务端可能不是同一版本。若消息结构只按当前代码隐式解析，一次滚动升级就可能造成字段错读。稳健协议会在 header 中放 protocol version、message type、flags、payload length 和 checksum。新增字段要有默认值，删除字段要经历兼容期。分布式系统的接口不是函数签名，而是长期存在的协议合同。

\begin{lstlisting}[language=C++,caption={消息头部的语义字段}]
struct WireHeader {
    std::uint16_t protocol_version = 1;
    std::uint16_t message_type = 0;
    std::uint32_t flags = 0;
    std::uint64_t request_id = 0;
    std::uint32_t payload_bytes = 0;
    std::uint32_t checksum = 0;
};
\end{lstlisting}

这段代码只是表达字段，不是跨平台二进制协议的完整实现。真实协议还要定义字节序、对齐、字段编码、字符串格式和兼容规则。不要直接把 C++ 结构体内存写到网络里，因为 padding、大小端、编译器布局和 ABI 都可能不同。协议格式要按字节定义，而不是按当前进程的内存布局定义。

\topic{序列化的成本和边界}

序列化不是简单的“把对象转成字节”。它会消耗 CPU、分配内存、访问字段、复制数据、计算校验，可能还要压缩。对象布局如果适合本地计算，不一定适合网络传输；网络格式如果紧凑，也不一定适合 CPU 直接计算。高性能系统常在内部布局和 wire format 之间做显式转换。

零拷贝也不是免费。要真正避免拷贝，需要缓冲区生命周期跨越发送完成，不能在发送返回后立即复用；需要处理分片、对齐和 scatter-gather；还要保证应用不能在内核或网卡尚未发送完成时修改内容。许多系统先用清晰的复制式协议做正确，再对热点路径引入缓冲池和批量发送。不要在协议语义没清楚时追求零拷贝。

序列化还影响故障恢复。检查点文件、shuffle manifest、任务提交表都需要稳定格式。若格式依赖当前类定义，代码升级后旧检查点可能无法读取。一个工程化系统应为持久化数据定义版本和迁移策略。分布式计算不是只在网络上传消息，还要把状态写到磁盘，并在失败后读回来。

\topic{逻辑时间和因果关系}

分布式系统不能依赖全局精确时钟来判断所有事件顺序。机器之间时钟会漂移，NTP 会校正，进程会暂停，网络延迟不稳定。物理时间适合做超时、租约和观测，但协议正确性常需要逻辑时间。最简单的逻辑时间是单调递增版本号、epoch、term 或日志索引。它们不表示真实时间，而表示协议中的先后和所有权。

例如 driver 每次发布新的路由表版本，就把 epoch 加一。Worker 收到旧 epoch 的任务，可以拒绝并要求刷新。Raft 的 term 也是类似思想：旧任期的 leader 即使还活着，也不能继续以旧身份提交新写。逻辑时间让系统能区分“我看到的状态”和“当前合法状态”。

向量时钟能表达更细的因果关系，但成本更高。很多系统不需要完整向量时钟，只需要针对某个资源的版本号。教材阶段要让读者理解：当两个事件分布在不同机器上时，肉眼时间先后不一定可靠；协议需要自己的顺序标记。

\topic{请求 ID 的设计}

请求 ID 不是随机数装饰，而是去重、追踪、幂等和调试的核心。一个请求 ID 至少应在一次逻辑操作内稳定，重试 attempt 不应改变逻辑 request id。否则服务端无法知道两次到达的请求其实是同一个操作的重复。Attempt id 可以单独增加，用于观测和超时分析。

请求 ID 的生成要考虑冲突范围。单机自增 ID 在单个进程内简单，但跨进程和重启后需要节点 ID、时间段或持久化计数。随机 ID 冲突概率低，但调试时不如结构化 ID 直观。雪花类 ID 把时间、节点和序列号组合起来，但依赖时钟行为。教学项目可以使用 driver 分配的 \code{std::uint64_t} 递增 ID，因为它在本机模拟里足够清楚；真实系统要写明冲突域和重启策略。

\begin{lstlisting}[language=C++,caption={逻辑请求和 attempt 分离}]
struct RequestKey {
    std::uint64_t request_id = 0;
    std::uint32_t attempt = 0;
};
\end{lstlisting}

服务端去重表通常按 \code{request_id} 存储逻辑结果，而不是按 attempt 存储。新 attempt 到达时，如果 \code{request_id} 已提交，就返回旧结果或忽略重复副作用。Attempt 只用于记录“这是第几次尝试”，不改变业务语义。

\topic{幂等表的容量和回收}

幂等表不能无限增长。服务端记录所有 request id，内存迟早耗尽。系统必须定义 request id 的有效期、客户端重试窗口和服务端保留时间。若客户端可能在一天后重试，服务端至少要保留一天；若服务端只保留一分钟，客户端一分钟后重试就可能重复执行。幂等语义是客户端和服务端共同的合同。

回收幂等表时要小心。按时间删除最简单，但依赖时钟和请求延迟分布。按提交日志截断更稳：如果所有客户端都保证不会重试某个 checkpoint 之前的请求，服务端可以删除旧记录。对于批处理任务，幂等表可以和 task attempt manifest 绑定：一个 stage 完成并 checkpoint 后，旧 attempt 去重记录可以随 checkpoint 压缩。

幂等表里保存什么也要权衡。只保存 request id 可以防止重复执行，但重复请求无法返回原响应；保存 request id 和响应结果，可以让客户端收到同样结果，但占用更多空间；保存结果摘要可以减少空间，但可能不能满足业务。设计时要从调用者语义出发。

\topic{错误分类}

分布式错误不能只返回“失败”。至少要区分可重试错误、不可重试错误、过载错误、版本错误、权限错误和未知错误。可重试错误包括临时网络失败、服务端 busy、leader 切换中。不可重试错误包括请求格式错误、权限不足、资源不存在。版本错误说明客户端路由表或协议版本过旧，需要刷新而不是盲目重试。

错误分类直接影响重试策略。对不可重试错误反复重试，只会浪费资源；对 busy 错误立即高频重试，会加剧过载；对版本错误不刷新路由，会一直打到错误节点。返回码是背压和控制面的接口，不只是调试信息。

\begin{lstlisting}[language=C++,caption={错误分类的教学形态}]
enum class RpcStatus : std::int32_t {
    Ok,
    Retryable,
    Busy,
    StaleRoute,
    InvalidRequest,
    PermissionDenied,
    Unknown,
};
\end{lstlisting}

真实协议还会有更细错误码和错误详情，但分类必须稳定。客户端重试库通常只看分类，日志和排查再看详细原因。把所有错误都归为 Unknown，会让客户端无法做正确行为。

\topic{超时预算的分解}

一个用户请求的总超时预算应被分解到各个下游调用，而不是每一层都重新给自己一个完整超时。假设用户请求总预算 1000 ms，服务 A 调服务 B，B 调 C。如果每层都设置 1000 ms，最坏情况下总延迟会远超用户预算。更好的做法是把剩余 deadline 随请求传递，下游根据剩余时间决定是否执行、降级或立即失败。

Deadline 比 timeout 更适合跨层传播。Timeout 是相对时间，每一层都要重新计算；deadline 是绝对截止点，传递后各层可以看当前剩余时间。时钟漂移会影响跨机器 deadline，因此通常还要保守处理，但它仍比无预算重试更清楚。批处理系统也有类似概念：一个 stage 的时间预算应来自整作业 SLA，而不是每个 task 自己无限重试。

\topic{重试、退避和抖动}

重试要慢下来，而不是立刻把同一个请求再打一遍。指数退避让连续失败后的等待时间增长，抖动让大量客户端不要在同一时刻重试。没有抖动时，所有客户端可能在服务恢复的一瞬间同时重试，造成第二波过载。重试预算限制一次逻辑请求最多产生多少 attempt，避免多层服务重试相乘。

重试还要和幂等绑定。读请求通常更容易重试，写请求必须有 request id 和提交语义。若写请求不是幂等的，客户端在超时后可能只能查询状态，而不是直接重试。教材要让读者看到：重试是正确性功能，不是简单可靠性增强。

\topic{慢节点和尾延迟}

分布式系统的整体延迟常被最慢分片决定。一个查询要访问十个分片，九个分片很快，一个分片慢，整体就慢。慢节点可能没有完全失败，只是磁盘抖动、CPU 被抢占、GC、网络拥塞或热点 key。故障检测如果只判断死活，就看不见慢节点。

处理慢节点的方法包括 speculative execution、请求 hedging、负载隔离、热点迁移和降级。Speculation 会启动重复 attempt，最先完成者获胜；hedging 会在等待一段时间后给另一个副本发同样读请求。这些方法能降低尾延迟，但增加资源消耗和重复副作用风险。没有幂等提交和重试预算，speculation 会让系统更危险。

\topic{队列和排队论直觉}

分布式系统里很多延迟来自排队。服务时间接近满载时，队列等待会快速上升。即使平均处理能力略高于平均到达率，突发流量也会造成队列。背压和限流的目的，就是在队列失控前让上游看到压力。

Little's Law 给出一个粗略关系：系统内平均在途请求数等于吞吐乘以平均延迟。若 in-flight 请求数持续增加而吞吐没有增加，说明延迟正在上升。若队列长度上升但 CPU 很空，可能是 I/O 或锁等待；若 CPU 满且队列上升，可能是计算能力不足；若网络出站队列上升，可能是下游慢或带宽瓶颈。队列是分布式系统最重要的压力表。

\topic{消息总线的本机实现}

本机模拟可以用一个 MessageBus 对象表示网络。它不需要真实 socket，先用多个有界队列表达节点之间的通道。每条消息带 source、target、request id、attempt、payload bytes、deadline 和 status。MessageBus 可以注入延迟、丢弃、重复和乱序。这样一台手机或一台开发机就能训练分布式协议。

\begin{lstlisting}[language=C++,caption={本机消息的核心字段}]
struct LocalMessage {
    std::uint32_t source_worker = 0;
    std::uint32_t target_worker = 0;
    std::uint64_t request_id = 0;
    std::uint32_t attempt = 0;
    std::uint32_t payload_bytes = 0;
    RpcStatus status = RpcStatus::Ok;
};
\end{lstlisting}

模拟器要把故障做成可重复。使用固定 seed，记录每次丢包、重复、延迟和重排事件。测试失败后，可以用同一个 seed 重放。没有可重复性，分布式 bug 会非常难查。

\topic{实验报告的故障矩阵}

分布式实验报告不能只写正常路径。至少要列出故障矩阵：请求丢失、响应丢失、服务端忙、服务端执行后崩溃、客户端超时后重试、重复消息到达、旧路由请求、慢节点、队列满、driver 重启。每一项都要写预期行为：重试、去重、返回旧结果、刷新路由、背压、失败还是恢复。

这张矩阵会迫使设计暴露缺口。例如如果响应丢失后客户端重试，而服务端没有保存响应结果，客户端可能无法知道第一次是否成功；如果执行后崩溃但提交点没持久化，恢复后可能重复执行；如果旧路由请求没有版本检查，迁移期间可能写到旧分片。故障矩阵不是测试附录，而是分布式设计本身。

\topic{连接管理}

真实 RPC 系统还要管理连接。每个请求新建 TCP 连接成本高，通常会复用连接池；连接池又要处理最大连接数、空闲超时、健康检查、连接重置和负载均衡。一个连接上的请求可能按顺序发送，也可能多路复用。若一个连接出现队头阻塞，后续请求会被拖慢；若连接太多，内核和远端都会有压力。

连接池的背压语义也要定义。连接都忙时，新请求是等待、失败、还是新建连接？等待是否消耗用户 deadline？连接错误时，正在飞行的请求如何标记？这些问题和线程池、有界队列非常相似。连接池不是网络库细节，而是分布式系统的执行资源管理。

\topic{服务发现和负载均衡}

客户端需要知道请求发给哪些服务实例。服务发现可以来自静态配置、DNS、注册中心或控制面。发现结果通常会缓存，因此要处理实例上下线、健康状态和版本变化。负载均衡策略可以是轮询、随机、最少连接、按延迟加权、按 locality 选择。策略不同，故障时行为也不同。

负载均衡不能只看请求数量。一个实例可能处理轻请求很多，另一个处理重请求少；一个实例延迟高可能是过载，也可能是网络远；一个实例刚启动缓存冷，直接大量流量打过去会抖动。工程系统常使用慢启动、健康检查、熔断和 outlier detection。教材项目可以简化为 driver 分配 worker，但应让读者知道真实服务发现是控制面的一部分。

\topic{熔断和隔离}

熔断器在下游持续失败或延迟过高时，暂时停止向它发送请求，快速失败或降级。这样可以避免线程、连接和队列被坏下游拖住，也给下游恢复时间。熔断不是为了隐藏错误，而是保护系统。熔断状态通常包括 closed、open、half-open：正常发送、快速失败、试探恢复。

隔离则是把不同下游或不同优先级请求放进不同资源池。若所有下游共享一个线程池，某个慢下游可能耗尽所有 worker，影响无关请求。隔离可以用独立队列、独立连接池、独立线程池或配额。代价是资源利用率可能下降，但故障影响范围更可控。

\topic{幂等性和业务键}

请求 ID 可以由系统生成，也可以来自业务键。创建订单时，客户端可以提供 client request id；写日志时，可以使用文件分片和 offset；提交 map output 时，可以使用 job id、stage id、task id。业务键更容易跨重启恢复，因为它和逻辑操作绑定。系统生成的随机 ID 如果只在内存里保存，重启后可能丢失去重语义。

幂等键还要定义作用域。一个 request id 在单个用户内唯一，还是全系统唯一？一个 task id 在一个 job 内唯一，还是跨 job 唯一？作用域不清会导致误去重或漏去重。教学项目应使用结构化 key，例如 job id、stage id、task id、attempt。这样日志和 manifest 更容易读。

\topic{时钟和单调时间}

超时测量应使用单调时钟，而不是 wall clock。系统时间可能被 NTP 调整，向前或向后跳；单调时钟适合计算 elapsed time。分布式系统传 deadline 时，跨机器 wall clock 差异仍要考虑，但本机队列和 I/O 超时应使用单调时间。日志时间戳和超时计时是不同用途。

租约和跨机器 deadline 更复杂。若依赖物理时间做安全判断，就要给时钟漂移和暂停留余量。很多协议避免直接依赖绝对时间来保证一致性，而使用 term、epoch 和 quorum。时间可以优化性能，不应轻易成为唯一正确性依据。

\topic{有序、因果和重放}

消息重放是测试分布式系统的重要方法。记录一段消息历史，包括发送、延迟、丢弃、重复和处理结果，然后在本机重放。若系统确定性足够，重放能复现 bug。为了支持重放，消息处理应尽量由输入消息和当前状态决定，避免隐藏随机数和真实时间直接影响核心语义。随机故障应使用记录的 seed。

因果关系也可以在 trace 中表达。一个 retry attempt 由哪次 timeout 触发，一个 reduce commit 由哪些 shuffle block 组成，一个 checkpoint 覆盖到哪个 stage，这些都应有 ID 关联。没有因果关联的日志，只是一堆时间戳，出了问题很难追。

\topic{分布式观测性}

分布式观测性至少包含 metrics、logs 和 traces。Metrics 聚合数量和延迟，适合看趋势；logs 记录离散事件和错误，适合排查细节；traces 串起一次请求或一个 task attempt 的路径，适合分析长尾。三者需要共同的 ID：request id、job id、stage id、task id、attempt、partition。没有 ID，无法把 driver、worker、shuffle 和 reduce 的事件连起来。

指标也要按维度分解。只看全局平均延迟，会掩盖某个 partition 长尾；只看总重试数，无法知道哪个 worker 或哪类错误导致；只看总吞吐，无法发现队列越来越长。分布式系统的指标设计要服务故障定位，而不是只服务漂亮仪表盘。

\topic{本章新增实验}

本章可以增加三个小实验。第一，连接池模拟：限制最大 in-flight 连接数，观察过载时等待和失败。第二，熔断模拟：让一个 worker 持续超时，客户端在多次失败后暂时停止发送，观察重试量是否下降。第三，trace 关联：为每条消息记录 request id 和 parent event id，重放一次响应丢失导致的重试。

这些实验不需要真实网络，但能训练分布式语义。它们也会给第 18 章的 driver 和 message bus 提供基础。

\topic{多租户和公平性}

分布式计算系统常同时运行多个 job。若所有 job 共用 worker 和 shuffle 队列，大 job 可能占满资源，小 job 长时间等待。公平调度需要定义资源份额、优先级和抢占策略。最简单的策略是每个 job 有独立队列，driver 按轮询或权重取任务；更复杂策略会考虑输入大小、deadline、用户配额和历史用量。

公平性和吞吐有时冲突。让大 job 连续占用资源，吞吐可能最高；穿插小 job，平均响应更好但缓存和调度成本增加。教学项目可以先支持单 job，但设计字段时保留 job id，就是为了未来多租户和公平调度。没有 job id，所有指标和去重都会混在一起。

\topic{背压和公平的交汇}

多个上游共享一个下游时，背压还要公平。若 reduce partition 队列满，是否阻塞所有 map worker，还是只阻塞发送该 partition 的 worker？若某个 job 产生大量热点 key，是否影响其他 job 的正常 partition？系统可以为每个 job 或 partition 设置独立队列和配额，防止一个热点拖垮全局。

这和单机线程池里的优先级队列类似。资源隔离降低故障传播，但可能降低资源利用率。工程设计要写清目标：批处理吞吐、交互延迟、公平性还是隔离。不同目标会得到不同调度器。

\topic{数据本地性和复制选择}

若输入分片有多个副本，调度器可以选择负载较低且距离较近的副本。只看本地性可能把任务都打到一个繁忙节点，只看负载可能产生大量远程读取。调度器需要一个成本函数，把数据位置、队列长度、节点健康、网络成本和优先级结合起来。教学项目可以用简单规则：优先本地，若本地队列超过阈值则远程。

这个规则和 NUMA 调度完全相似：优先本地内存，过载时跨节点。第二册反复建立这种跨尺度类比，是为了让读者面对新系统时能迁移思维。

\topic{协议演进和兼容测试}

分布式协议一旦持久化或跨进程，就需要兼容测试。新增字段要有默认值，旧节点收到新字段能忽略，新节点收到旧消息能补默认值。删除字段要经历过渡期。消息版本、feature flags 和最小兼容版本都要有策略。否则滚动升级时，集群可能出现部分新旧节点互相不能通信。

教学项目可以简单要求“同一版本运行”，但如果写成教材，就要说明这是简化。第三册 AI 推理引擎如果加入模型格式和算子版本，也会遇到同样问题。协议演进是长期系统的基本能力。

\topic{混沌测试和演练}

故障注入可以从单元测试扩展到混沌测试：在长时间运行中随机让 worker 变慢、丢消息、重启 driver、损坏临时 block、让队列变满。混沌测试的目标不是制造混乱，而是验证系统在预期故障集合内保持可控。每次注入都应记录事件，失败后能回放。

混沌测试必须有安全边界。不要在没有隔离的真实数据上随意注入；不要注入系统无法承受的未知破坏；不要只看系统没崩就算通过，还要检查结果正确、重试预算、队列长度和恢复时间。可靠性测试同样需要 oracle。

\topic{分布式模型的知识地图}

到本章结束，读者应能把分布式系统拆成：消息、时间、故障、身份、路由、幂等、背压、调度和观测。消息回答数据如何传；时间回答何时放弃；故障回答什么可能坏；身份回答这是哪个请求；路由回答发给谁；幂等回答重复怎么办；背压回答过载怎么办；调度回答谁执行；观测回答如何知道发生了什么。后两章会把这些部件组合成分片复制和计算引擎。

\topic{分布式模型章节总结}

分布式系统的困难不是网络 API，而是不确定性。消息可能丢、重复、乱序；节点可能慢、停、恢复；客户端超时不代表服务端停止；重试可能放大故障；旧路由和旧 leader 可能继续发送请求。解决这些问题需要请求身份、版本、deadline、幂等表、背压、故障矩阵和 trace。把这些基础模型建好，后面的复制和计算工程才有可靠语义。

\topic{分布式模型决策表}

设计一个 RPC 或消息协议时，先填决策表。请求是否幂等？若不是，超时后不能盲目重试。是否有 request id？若没有，服务端无法去重。是否有 deadline？若没有，多层调用会放大延迟。是否区分 busy 和 fatal？若没有，客户端无法正确退避。是否有路由版本？若没有，迁移期间旧请求难以拒绝。是否有 trace id？若没有，故障后无法串联路径。

这张表适用于任何分布式调用，也适用于本机异步消息。只要执行和完成分离，就需要身份、时间、状态和错误语义。

\topic{消息协议审查表}

每个消息类型都应审查：谁发送，谁接收，是否允许重复，是否有响应，超时后发送方做什么，接收方是否幂等，payload 是否有版本和长度，错误码是否分类，是否进入 trace，是否影响提交状态。审查后会发现很多消息只是指标，可以丢；有些消息是提交，必须持久化；有些消息是控制面，必须带 epoch。消息不是字节包，而是系统状态转换。


\topic{高密度主线补充：从失败推导机制}

分布式模型章要从一次函数调用变成一次消息后丢失了什么保证讲起。这一轮补充专门解决两个问题：第一，避免把本章写成概念枚举；第二，让每个概念都能从driver 向 worker 发送任务并等待结果的失败中自然推出。本章的核心失败不是“代码不会写”，而是远端执行的完成、失败和重复都无法靠本地调用栈表达。所以正文的顺序必须始终保持为：先给出具体任务，再写朴素方案，再观察失败信号，再引入机制，最后给出实验和边界。
这样的写法比直接给定义慢一些，但它更接近工程学习的真实路径。真实系统很少把问题包装成选择题；它只会表现为吞吐上不去、CPU 利用率怪、结果偶尔错、队列越积越多、重启后状态对不上、线上延迟突然变长。读者要学会从这些现象反推机制，而不是看到术语才想起术语。

\topic{消息身份}

先把问题收窄到一个可推导的场景：响应回来时如何知道它属于哪个请求和哪个 attempt。在driver 向 worker 发送任务并等待结果里，这个问题不会以术语形式出现，而是以结果不稳定、吞吐不增长、延迟抖动或恢复困难的形式出现。如果一上来只背消息身份的定义，读者很容易把它当成孤立知识点；更好的顺序是先看朴素方案为什么合理，再看它在真实约束下怎样失败。

朴素方案通常是：只按连接或 worker 匹配响应。这个方案的吸引力在于它贴近单线程直觉，代码短，状态少，测试也容易写。但是它隐含了几个没有说出口的假设：输入总是干净，资源近似无限，执行不会交错，失败不会发生，硬件成本和源码结构大体一致。一旦这些假设被打破，重试、乱序和重复响应会污染状态。这时继续在原方案上加 if 或加锁，往往只会把真正的问题埋得更深。

消息身份就是从这个失败里长出来的概念。它的核心模型可以这样理解：request id、attempt 和 epoch 让消息具有可验证身份。这个模型必须同时放在三个层次看。源码层要说明谁读谁写、谁拥有对象、哪个状态允许变化；运行时层要说明线程、队列、任务和 I/O 怎样排队；硬件或系统层要说明缓存、页、调度、文件系统或网络怎样参与成本。三层缺一层，解释都会变形：只有源码层会看不见隐藏成本，只有硬件层会丢失业务语义，只有运行时层又容易把正确性和性能混在一起。

观察方法不能停在“跑一下变快了”。这一点在本章尤其重要：在模拟网络中打乱响应顺序。实验要有 reference，要固定输入规模，要记录环境，要把冷启动、预热、核心循环、提交和清理分开。如果工具权限不足，也要写清楚不能观察什么，而不是把猜测写成结论。系统教材的严谨性来自证据链：现象、假设、对照、指标、结论边界。

消息身份也有边界。身份字段必须进入提交边界，不只是日志字段。工程判断不是把一个技术推到所有地方，而是知道它在哪些约束下成立。读者在本章应形成一种习惯：每学到一个机制，都要问它解决了什么失败、引入了什么新成本、需要什么测试、在什么输入或部署环境下会失效。

\topic{Deadline}

先把问题收窄到一个可推导的场景：超时应该从每一层重新计算还是继承调用预算。在driver 向 worker 发送任务并等待结果里，这个问题不会以术语形式出现，而是以结果不稳定、吞吐不增长、延迟抖动或恢复困难的形式出现。如果一上来只背Deadline的定义，读者很容易把它当成孤立知识点；更好的顺序是先看朴素方案为什么合理，再看它在真实约束下怎样失败。

朴素方案通常是：每层都设置固定超时。这个方案的吸引力在于它贴近单线程直觉，代码短，状态少，测试也容易写。但是它隐含了几个没有说出口的假设：输入总是干净，资源近似无限，执行不会交错，失败不会发生，硬件成本和源码结构大体一致。一旦这些假设被打破，多层调用叠加后总时间失控。这时继续在原方案上加 if 或加锁，往往只会把真正的问题埋得更深。

Deadline就是从这个失败里长出来的概念。它的核心模型可以这样理解：deadline 表示绝对预算，沿调用链传递。这个模型必须同时放在三个层次看。源码层要说明谁读谁写、谁拥有对象、哪个状态允许变化；运行时层要说明线程、队列、任务和 I/O 怎样排队；硬件或系统层要说明缓存、页、调度、文件系统或网络怎样参与成本。三层缺一层，解释都会变形：只有源码层会看不见隐藏成本，只有硬件层会丢失业务语义，只有运行时层又容易把正确性和性能混在一起。

观察方法不能停在“跑一下变快了”。这一点在本章尤其重要：记录剩余时间和超时原因。实验要有 reference，要固定输入规模，要记录环境，要把冷启动、预热、核心循环、提交和清理分开。如果工具权限不足，也要写清楚不能观察什么，而不是把猜测写成结论。系统教材的严谨性来自证据链：现象、假设、对照、指标、结论边界。

Deadline也有边界。过短 deadline 会造成无意义重试，过长会占住资源。工程判断不是把一个技术推到所有地方，而是知道它在哪些约束下成立。读者在本章应形成一种习惯：每学到一个机制，都要问它解决了什么失败、引入了什么新成本、需要什么测试、在什么输入或部署环境下会失效。

\topic{重试}

先把问题收窄到一个可推导的场景：请求超时后是否可以直接再发一次。在driver 向 worker 发送任务并等待结果里，这个问题不会以术语形式出现，而是以结果不稳定、吞吐不增长、延迟抖动或恢复困难的形式出现。如果一上来只背重试的定义，读者很容易把它当成孤立知识点；更好的顺序是先看朴素方案为什么合理，再看它在真实约束下怎样失败。

朴素方案通常是：立即重试直到成功。这个方案的吸引力在于它贴近单线程直觉，代码短，状态少，测试也容易写。但是它隐含了几个没有说出口的假设：输入总是干净，资源近似无限，执行不会交错，失败不会发生，硬件成本和源码结构大体一致。一旦这些假设被打破，服务端可能已经执行，重复请求可能重复提交。这时继续在原方案上加 if 或加锁，往往只会把真正的问题埋得更深。

重试就是从这个失败里长出来的概念。它的核心模型可以这样理解：重试需要预算、退避、抖动和幂等语义。这个模型必须同时放在三个层次看。源码层要说明谁读谁写、谁拥有对象、哪个状态允许变化；运行时层要说明线程、队列、任务和 I/O 怎样排队；硬件或系统层要说明缓存、页、调度、文件系统或网络怎样参与成本。三层缺一层，解释都会变形：只有源码层会看不见隐藏成本，只有硬件层会丢失业务语义，只有运行时层又容易把正确性和性能混在一起。

观察方法不能停在“跑一下变快了”。这一点在本章尤其重要：注入响应丢失和服务端慢处理。实验要有 reference，要固定输入规模，要记录环境，要把冷启动、预热、核心循环、提交和清理分开。如果工具权限不足，也要写清楚不能观察什么，而不是把猜测写成结论。系统教材的严谨性来自证据链：现象、假设、对照、指标、结论边界。

重试也有边界。不可幂等操作不能盲目重试。工程判断不是把一个技术推到所有地方，而是知道它在哪些约束下成立。读者在本章应形成一种习惯：每学到一个机制，都要问它解决了什么失败、引入了什么新成本、需要什么测试、在什么输入或部署环境下会失效。

\topic{幂等}

先把问题收窄到一个可推导的场景：重复执行同一个请求为什么不应改变最终结果。在driver 向 worker 发送任务并等待结果里，这个问题不会以术语形式出现，而是以结果不稳定、吞吐不增长、延迟抖动或恢复困难的形式出现。如果一上来只背幂等的定义，读者很容易把它当成孤立知识点；更好的顺序是先看朴素方案为什么合理，再看它在真实约束下怎样失败。

朴素方案通常是：服务端每收到一次就执行一次。这个方案的吸引力在于它贴近单线程直觉，代码短，状态少，测试也容易写。但是它隐含了几个没有说出口的假设：输入总是干净，资源近似无限，执行不会交错，失败不会发生，硬件成本和源码结构大体一致。一旦这些假设被打破，响应丢失后重试会重复计数或重复写输出。这时继续在原方案上加 if 或加锁，往往只会把真正的问题埋得更深。

幂等就是从这个失败里长出来的概念。它的核心模型可以这样理解：幂等表记录 request id、状态和结果，使重复请求返回同一结果。这个模型必须同时放在三个层次看。源码层要说明谁读谁写、谁拥有对象、哪个状态允许变化；运行时层要说明线程、队列、任务和 I/O 怎样排队；硬件或系统层要说明缓存、页、调度、文件系统或网络怎样参与成本。三层缺一层，解释都会变形：只有源码层会看不见隐藏成本，只有硬件层会丢失业务语义，只有运行时层又容易把正确性和性能混在一起。

观察方法不能停在“跑一下变快了”。这一点在本章尤其重要：测试执行成功但响应丢失。实验要有 reference，要固定输入规模，要记录环境，要把冷启动、预热、核心循环、提交和清理分开。如果工具权限不足，也要写清楚不能观察什么，而不是把猜测写成结论。系统教材的严谨性来自证据链：现象、假设、对照、指标、结论边界。

幂等也有边界。幂等记录本身也要和业务提交同边界。工程判断不是把一个技术推到所有地方，而是知道它在哪些约束下成立。读者在本章应形成一种习惯：每学到一个机制，都要问它解决了什么失败、引入了什么新成本、需要什么测试、在什么输入或部署环境下会失效。

\topic{部分失败}

先把问题收窄到一个可推导的场景：客户端活着、服务端活着但网络断开时系统处于什么状态。在driver 向 worker 发送任务并等待结果里，这个问题不会以术语形式出现，而是以结果不稳定、吞吐不增长、延迟抖动或恢复困难的形式出现。如果一上来只背部分失败的定义，读者很容易把它当成孤立知识点；更好的顺序是先看朴素方案为什么合理，再看它在真实约束下怎样失败。

朴素方案通常是：把失败视为进程崩溃或成功两种。这个方案的吸引力在于它贴近单线程直觉，代码短，状态少，测试也容易写。但是它隐含了几个没有说出口的假设：输入总是干净，资源近似无限，执行不会交错，失败不会发生，硬件成本和源码结构大体一致。一旦这些假设被打破，真实系统存在消息丢失、延迟、分区和旧 leader。这时继续在原方案上加 if 或加锁，往往只会把真正的问题埋得更深。

部分失败就是从这个失败里长出来的概念。它的核心模型可以这样理解：部分失败意味着一方无法可靠知道另一方状态。这个模型必须同时放在三个层次看。源码层要说明谁读谁写、谁拥有对象、哪个状态允许变化；运行时层要说明线程、队列、任务和 I/O 怎样排队；硬件或系统层要说明缓存、页、调度、文件系统或网络怎样参与成本。三层缺一层，解释都会变形：只有源码层会看不见隐藏成本，只有硬件层会丢失业务语义，只有运行时层又容易把正确性和性能混在一起。

观察方法不能停在“跑一下变快了”。这一点在本章尤其重要：用消息丢弃、重复、延迟和乱序模拟。实验要有 reference，要固定输入规模，要记录环境，要把冷启动、预热、核心循环、提交和清理分开。如果工具权限不足，也要写清楚不能观察什么，而不是把猜测写成结论。系统教材的严谨性来自证据链：现象、假设、对照、指标、结论边界。

部分失败也有边界。超时只是未知，不是远端已停止。工程判断不是把一个技术推到所有地方，而是知道它在哪些约束下成立。读者在本章应形成一种习惯：每学到一个机制，都要问它解决了什么失败、引入了什么新成本、需要什么测试、在什么输入或部署环境下会失效。

\topic{流控}

先把问题收窄到一个可推导的场景：worker 返回 busy 时 driver 应该怎么做。在driver 向 worker 发送任务并等待结果里，这个问题不会以术语形式出现，而是以结果不稳定、吞吐不增长、延迟抖动或恢复困难的形式出现。如果一上来只背流控的定义，读者很容易把它当成孤立知识点；更好的顺序是先看朴素方案为什么合理，再看它在真实约束下怎样失败。

朴素方案通常是：继续提交任务直到队列满。这个方案的吸引力在于它贴近单线程直觉，代码短，状态少，测试也容易写。但是它隐含了几个没有说出口的假设：输入总是干净，资源近似无限，执行不会交错，失败不会发生，硬件成本和源码结构大体一致。一旦这些假设被打破，重试和排队放大故障。这时继续在原方案上加 if 或加锁，往往只会把真正的问题埋得更深。

流控就是从这个失败里长出来的概念。它的核心模型可以这样理解：流控把接收方容量反馈给发送方。这个模型必须同时放在三个层次看。源码层要说明谁读谁写、谁拥有对象、哪个状态允许变化；运行时层要说明线程、队列、任务和 I/O 怎样排队；硬件或系统层要说明缓存、页、调度、文件系统或网络怎样参与成本。三层缺一层，解释都会变形：只有源码层会看不见隐藏成本，只有硬件层会丢失业务语义，只有运行时层又容易把正确性和性能混在一起。

观察方法不能停在“跑一下变快了”。这一点在本章尤其重要：记录 in flight、busy 响应和退避时间。实验要有 reference，要固定输入规模，要记录环境，要把冷启动、预热、核心循环、提交和清理分开。如果工具权限不足，也要写清楚不能观察什么，而不是把猜测写成结论。系统教材的严谨性来自证据链：现象、假设、对照、指标、结论边界。

流控也有边界。流控不是错误，它是保护系统的协议。工程判断不是把一个技术推到所有地方，而是知道它在哪些约束下成立。读者在本章应形成一种习惯：每学到一个机制，都要问它解决了什么失败、引入了什么新成本、需要什么测试、在什么输入或部署环境下会失效。

\topic{Trace Id}

先把问题收窄到一个可推导的场景：一次作业跨多个 worker 后如何定位慢在哪里。在driver 向 worker 发送任务并等待结果里，这个问题不会以术语形式出现，而是以结果不稳定、吞吐不增长、延迟抖动或恢复困难的形式出现。如果一上来只背Trace Id的定义，读者很容易把它当成孤立知识点；更好的顺序是先看朴素方案为什么合理，再看它在真实约束下怎样失败。

朴素方案通常是：每个进程各写自己的日志。这个方案的吸引力在于它贴近单线程直觉，代码短，状态少，测试也容易写。但是它隐含了几个没有说出口的假设：输入总是干净，资源近似无限，执行不会交错，失败不会发生，硬件成本和源码结构大体一致。一旦这些假设被打破，无法串联同一请求的路径。这时继续在原方案上加 if 或加锁，往往只会把真正的问题埋得更深。

Trace Id就是从这个失败里长出来的概念。它的核心模型可以这样理解：trace id 把跨组件事件连成一条时间线。这个模型必须同时放在三个层次看。源码层要说明谁读谁写、谁拥有对象、哪个状态允许变化；运行时层要说明线程、队列、任务和 I/O 怎样排队；硬件或系统层要说明缓存、页、调度、文件系统或网络怎样参与成本。三层缺一层，解释都会变形：只有源码层会看不见隐藏成本，只有硬件层会丢失业务语义，只有运行时层又容易把正确性和性能混在一起。

观察方法不能停在“跑一下变快了”。这一点在本章尤其重要：记录 driver send、worker receive、execute、reply。实验要有 reference，要固定输入规模，要记录环境，要把冷启动、预热、核心循环、提交和清理分开。如果工具权限不足，也要写清楚不能观察什么，而不是把猜测写成结论。系统教材的严谨性来自证据链：现象、假设、对照、指标、结论边界。

Trace Id也有边界。trace 也要采样和限流，不能让观测拖垮系统。工程判断不是把一个技术推到所有地方，而是知道它在哪些约束下成立。读者在本章应形成一种习惯：每学到一个机制，都要问它解决了什么失败、引入了什么新成本、需要什么测试、在什么输入或部署环境下会失效。

\topic{本机模拟}

先把问题收窄到一个可推导的场景：没有集群时怎样学习分布式失败。在driver 向 worker 发送任务并等待结果里，这个问题不会以术语形式出现，而是以结果不稳定、吞吐不增长、延迟抖动或恢复困难的形式出现。如果一上来只背本机模拟的定义，读者很容易把它当成孤立知识点；更好的顺序是先看朴素方案为什么合理，再看它在真实约束下怎样失败。

朴素方案通常是：只写正常路径函数调用。这个方案的吸引力在于它贴近单线程直觉，代码短，状态少，测试也容易写。但是它隐含了几个没有说出口的假设：输入总是干净，资源近似无限，执行不会交错，失败不会发生，硬件成本和源码结构大体一致。一旦这些假设被打破，无法暴露乱序、重复和超时。这时继续在原方案上加 if 或加锁，往往只会把真正的问题埋得更深。

本机模拟就是从这个失败里长出来的概念。它的核心模型可以这样理解：本机 message bus 可以注入延迟、丢失、重复和乱序。这个模型必须同时放在三个层次看。源码层要说明谁读谁写、谁拥有对象、哪个状态允许变化；运行时层要说明线程、队列、任务和 I/O 怎样排队；硬件或系统层要说明缓存、页、调度、文件系统或网络怎样参与成本。三层缺一层，解释都会变形：只有源码层会看不见隐藏成本，只有硬件层会丢失业务语义，只有运行时层又容易把正确性和性能混在一起。

观察方法不能停在“跑一下变快了”。这一点在本章尤其重要：用固定随机种子复现实验。实验要有 reference，要固定输入规模，要记录环境，要把冷启动、预热、核心循环、提交和清理分开。如果工具权限不足，也要写清楚不能观察什么，而不是把猜测写成结论。系统教材的严谨性来自证据链：现象、假设、对照、指标、结论边界。

本机模拟也有边界。模拟不能证明真实网络性能，但能训练协议语义。工程判断不是把一个技术推到所有地方，而是知道它在哪些约束下成立。读者在本章应形成一种习惯：每学到一个机制，都要问它解决了什么失败、引入了什么新成本、需要什么测试、在什么输入或部署环境下会失效。

\topic{案例与实验串联}

案例“响应丢失”用于把抽象机制落到一段可复现实验。场景是：worker 已经完成任务并提交结果，但回复在路上丢失。这个场景要足够小，小到读者可以手工画出数据流、状态变化和成本路径；又要足够真实，能暴露教材正文讨论的失败模式。

第一版不要急着写最终答案，而应保留一个朴素版本：driver 超时后派发新 attempt。朴素版本的价值不是性能，而是语义清楚。它告诉我们结果应该是什么，也告诉我们优化前系统在哪些地方偷懒。

第二版再引入改进：worker 或 reduce 端用 request id 去重。改进必须说明成本迁移到哪里，而不是只说“更快”。例如减少共享写可能增加最终合并成本，批处理可能增加尾延迟，分片可能引入负载倾斜，异步可能让生命周期更难验证。

验证时重点看：结果只提交一次，重复响应返回同一摘要。报告里至少写出输入规模、硬件或系统环境、编译选项、运行命令、关键指标和反例。若结果与预期不同，优先检查实验变量，而不是马上改代码。
案例“重试风暴”用于把抽象机制落到一段可复现实验。场景是：服务端变慢，客户端同时超时并重试。这个场景要足够小，小到读者可以手工画出数据流、状态变化和成本路径；又要足够真实，能暴露教材正文讨论的失败模式。

第一版不要急着写最终答案，而应保留一个朴素版本：所有客户端立即重试。朴素版本的价值不是性能，而是语义清楚。它告诉我们结果应该是什么，也告诉我们优化前系统在哪些地方偷懒。

第二版再引入改进：指数退避、抖动和全局 in flight 限制。改进必须说明成本迁移到哪里，而不是只说“更快”。例如减少共享写可能增加最终合并成本，批处理可能增加尾延迟，分片可能引入负载倾斜，异步可能让生命周期更难验证。

验证时重点看：请求量不会因重试超过预算。报告里至少写出输入规模、硬件或系统环境、编译选项、运行命令、关键指标和反例。若结果与预期不同，优先检查实验变量，而不是马上改代码。
案例“乱序响应”用于把抽象机制落到一段可复现实验。场景是：attempt 2 的响应先到，attempt 1 的响应后到。这个场景要足够小，小到读者可以手工画出数据流、状态变化和成本路径；又要足够真实，能暴露教材正文讨论的失败模式。

第一版不要急着写最终答案，而应保留一个朴素版本：谁后到就覆盖状态。朴素版本的价值不是性能，而是语义清楚。它告诉我们结果应该是什么，也告诉我们优化前系统在哪些地方偷懒。

第二版再引入改进：用 attempt 和 epoch 做状态检查。改进必须说明成本迁移到哪里，而不是只说“更快”。例如减少共享写可能增加最终合并成本，批处理可能增加尾延迟，分片可能引入负载倾斜，异步可能让生命周期更难验证。

验证时重点看：旧 attempt 不会覆盖已提交结果。报告里至少写出输入规模、硬件或系统环境、编译选项、运行命令、关键指标和反例。若结果与预期不同，优先检查实验变量，而不是马上改代码。

\topic{Linux 源码阅读与系统观察}

Linux 源码阅读要从用户态现象倒推入口。本章可围绕 \filepath{net/core/dev.c}、\filepath{net/ipv4/tcp.c}、\filepath{kernel/time/timer.c} 做窄范围阅读。不要试图一次读完整子系统，先选择一个问题，例如一次阻塞为什么睡眠、一次缺页怎样建立映射、一次唤醒怎样进入 run queue、一次写回怎样变成持久化边界。

阅读报告应固定内核版本、阅读路径和实验命令。第一段写用户态最小程序，第二段写观察到的现象，第三段写源码中的关键结构和状态变化，第四段写结论边界。若当前设备不是对应内核，报告必须说明源码只是机制参考，而不是当前运行系统的直接证据。

源码阅读还要看数据结构，而不只是函数名。队列、树、链表、位图、引用计数、状态枚举、等待队列、页表、inode、bio、task 结构这些细节，决定了机制怎样落地。读者要练习把一个用户态概念翻译成内核对象：线程对应 task，等待对应 wait queue 或 futex 路径，文件缓存对应 page cache，内存策略对应 VMA 和页面分配。

\topic{贯穿项目检查点}

把本章落到贯穿项目时，不要只加一个接口或一个 benchmark。每次新增能力都应有语义目标、最小实现、对照实验、指标和失败注入。项目不是堆功能，而是把本章的机制变成可运行、可检查、可复盘的工程对象。

\begin{lstlisting}[numbers=none]
chapter project checkpoints:
  1. 为每条消息加入 request id 和 attempt
  2. 实现可注入故障的 message bus
  3. 加入 deadline 和重试预算
  4. 实现幂等结果表
  5. 输出跨组件 trace
\end{lstlisting}

项目报告要回答四个问题。第一，朴素版本是什么，为什么不足。第二，改进版本改变了哪个边界，是数据边界、执行边界、同步边界、I/O 边界还是提交边界。第三，证据是什么，哪些指标支持结论。第四，剩余风险是什么，在哪些输入、机器或失败条件下结论可能不成立。只有这样，项目才不会变成代码展示，而会变成系统能力训练。



\topic{第二轮深水区：把机制讲成可验证能力}

本章继续扩写时，重点不再是补充更多名词，而是把已经出现的机制变成可验证能力。围绕driver 向 worker 发送任务并等待结果，读者最终应能做三件事：解释为什么朴素方案失败，设计能隔离变量的实验，把实验结论落到代码边界。这三件事缺一不可。只会解释会变成空谈，只会跑实验会变成工具使用，只会改代码又会在复杂系统里丢掉语义。
本章的主失败是：远端执行的完成、失败和重复都无法靠本地调用栈表达。这个失败会在不同尺度反复出现。小输入里它可能只是一次结果不稳定；多线程里它可能变成锁竞争、乱序或重复写；I/O 和分布式里它可能变成提交不清、恢复不可靠或重试放大。所以本章的每个机制都要写出尺度变化后的表现。

\topic{深挖：消息身份}

围绕“消息身份”，先把它放回一个具体问题：响应回来时如何知道它属于哪个请求和哪个 attempt。如果这个问题只在单线程小输入里出现，读者很容易把它当成局部技巧；但在driver 向 worker 发送任务并等待结果中，它会沿着输入、执行、同步、I/O 和提交一路传播。所以讲解时不能只写定义，而要让读者看到它怎样从一个小失败变成系统性约束。

朴素写法通常是“只按连接或 worker 匹配响应”。这句话看似简单，却包含几个默认前提：状态只有一个拥有者，执行顺序可预测，资源足够近，失败不会穿过边界，观察结果能够代表真实成本。这些前提在教学小程序中可能成立，在driver 向 worker 发送任务并等待结果里却会被输入规模、线程交错、硬件拓扑或故障注入逐个打破。

失败信号是“重试、乱序和重复响应会污染状态”。注意，这个失败不一定表现为崩溃。它也可能表现为吞吐不再增长、p99 抖动、重启后结果多了一份、某个分区长尾、CPU 使用率很高但有效工作很少。教材要把这些信号和机制绑定起来，否则读者只会背术语，遇到真实现象时仍然不知道从哪里下手。

更可靠的推导顺序是先写出不变量。对消息身份来说，不变量不是一句口号，而是本章要保护的事实：哪些数据只能写一次，哪些状态可以重试，哪些结果可以被缓存，哪些成本必须摊销，哪些错误必须外显。只要不变量没有写清，后面的代码、实验和优化都没有稳定坐标。

然后才引入模型：request id、attempt 和 epoch 让消息具有可验证身份。这个模型应同时解释正确性和成本。正确性部分回答“为什么结果可信”，成本部分回答“为什么这个实现会快或慢”。如果一个模型只能解释速度，却解释不了失败恢复，它不适合系统教材；如果只能解释语义，却完全不关心 cache、调度、I/O 或网络成本，也不适合第二册。

实验应围绕“只改变一个变量”设计。针对消息身份，最小实验可以保留朴素版本，再实现一个改进版本，输入规模从小到大，线程数或并发度从一到目标机器上限，错误注入从无到有。每一组实验都应记录 reference 结果、核心指标、环境和命令。这样读者看到差异时，可以判断差异来自机制本身，而不是来自初始化、缓存冷热、调度迁移或文件系统状态。

观察手段是：在模拟网络中打乱响应顺序。观察不是为了堆工具名，而是为了回答一个具体假设。若假设是共享写导致扩展性下降，就应看线程数曲线、写入布局和 cache 相关指标；若假设是提交边界错误，就应做崩溃点矩阵；若假设是队列背压，就应看水位、等待和上游速率。工具输出必须回到假设，不要把截图或命令输出当成结论。

最后要讲边界：身份字段必须进入提交边界，不只是日志字段。高质量教材必须把反例放在正文里。读者需要知道什么时候该用这个机制，什么时候它会制造新问题，什么时候应该退回更简单方案。比如一个单线程工具没有并发共享，强行引入复杂运行时只会增加维护成本；一个没有恢复要求的临时分析脚本，也不必承担完整 manifest 协议。

把这一点落实到代码审查，可以要求每个涉及消息身份的改动都回答五个问题：朴素版本是什么，新增机制保护了哪个不变量，新增成本在哪里，如何回退，哪些测试证明边界。这五个问题比“用了某某技术”更能判断质量，也能防止团队在没有证据时追逐复杂实现。

\topic{深挖：Deadline}

围绕“Deadline”，先把它放回一个具体问题：超时应该从每一层重新计算还是继承调用预算。如果这个问题只在单机多线程里出现，读者很容易把它当成局部技巧；但在driver 向 worker 发送任务并等待结果中，它会沿着输入、执行、同步、I/O 和提交一路传播。所以讲解时不能只写定义，而要让读者看到它怎样从一个小失败变成系统性约束。

朴素写法通常是“每层都设置固定超时”。这句话看似简单，却包含几个默认前提：状态只有一个拥有者，执行顺序可预测，资源足够近，失败不会穿过边界，观察结果能够代表真实成本。这些前提在教学小程序中可能成立，在driver 向 worker 发送任务并等待结果里却会被输入规模、线程交错、硬件拓扑或故障注入逐个打破。

失败信号是“多层调用叠加后总时间失控”。注意，这个失败不一定表现为崩溃。它也可能表现为吞吐不再增长、p99 抖动、重启后结果多了一份、某个分区长尾、CPU 使用率很高但有效工作很少。教材要把这些信号和机制绑定起来，否则读者只会背术语，遇到真实现象时仍然不知道从哪里下手。

更可靠的推导顺序是先写出不变量。对Deadline来说，不变量不是一句口号，而是本章要保护的事实：哪些数据只能写一次，哪些状态可以重试，哪些结果可以被缓存，哪些成本必须摊销，哪些错误必须外显。只要不变量没有写清，后面的代码、实验和优化都没有稳定坐标。

然后才引入模型：deadline 表示绝对预算，沿调用链传递。这个模型应同时解释正确性和成本。正确性部分回答“为什么结果可信”，成本部分回答“为什么这个实现会快或慢”。如果一个模型只能解释速度，却解释不了失败恢复，它不适合系统教材；如果只能解释语义，却完全不关心 cache、调度、I/O 或网络成本，也不适合第二册。

实验应围绕“只改变一个变量”设计。针对Deadline，最小实验可以保留朴素版本，再实现一个改进版本，输入规模从小到大，线程数或并发度从一到目标机器上限，错误注入从无到有。每一组实验都应记录 reference 结果、核心指标、环境和命令。这样读者看到差异时，可以判断差异来自机制本身，而不是来自初始化、缓存冷热、调度迁移或文件系统状态。

观察手段是：记录剩余时间和超时原因。观察不是为了堆工具名，而是为了回答一个具体假设。若假设是共享写导致扩展性下降，就应看线程数曲线、写入布局和 cache 相关指标；若假设是提交边界错误，就应做崩溃点矩阵；若假设是队列背压，就应看水位、等待和上游速率。工具输出必须回到假设，不要把截图或命令输出当成结论。

最后要讲边界：过短 deadline 会造成无意义重试，过长会占住资源。高质量教材必须把反例放在正文里。读者需要知道什么时候该用这个机制，什么时候它会制造新问题，什么时候应该退回更简单方案。比如一个单线程工具没有并发共享，强行引入复杂运行时只会增加维护成本；一个没有恢复要求的临时分析脚本，也不必承担完整 manifest 协议。

把这一点落实到代码审查，可以要求每个涉及Deadline的改动都回答五个问题：朴素版本是什么，新增机制保护了哪个不变量，新增成本在哪里，如何回退，哪些测试证明边界。这五个问题比“用了某某技术”更能判断质量，也能防止团队在没有证据时追逐复杂实现。

\topic{深挖：重试}

围绕“重试”，先把它放回一个具体问题：请求超时后是否可以直接再发一次。如果这个问题只在NUMA 或异构核心里出现，读者很容易把它当成局部技巧；但在driver 向 worker 发送任务并等待结果中，它会沿着输入、执行、同步、I/O 和提交一路传播。所以讲解时不能只写定义，而要让读者看到它怎样从一个小失败变成系统性约束。

朴素写法通常是“立即重试直到成功”。这句话看似简单，却包含几个默认前提：状态只有一个拥有者，执行顺序可预测，资源足够近，失败不会穿过边界，观察结果能够代表真实成本。这些前提在教学小程序中可能成立，在driver 向 worker 发送任务并等待结果里却会被输入规模、线程交错、硬件拓扑或故障注入逐个打破。

失败信号是“服务端可能已经执行，重复请求可能重复提交”。注意，这个失败不一定表现为崩溃。它也可能表现为吞吐不再增长、p99 抖动、重启后结果多了一份、某个分区长尾、CPU 使用率很高但有效工作很少。教材要把这些信号和机制绑定起来，否则读者只会背术语，遇到真实现象时仍然不知道从哪里下手。

更可靠的推导顺序是先写出不变量。对重试来说，不变量不是一句口号，而是本章要保护的事实：哪些数据只能写一次，哪些状态可以重试，哪些结果可以被缓存，哪些成本必须摊销，哪些错误必须外显。只要不变量没有写清，后面的代码、实验和优化都没有稳定坐标。

然后才引入模型：重试需要预算、退避、抖动和幂等语义。这个模型应同时解释正确性和成本。正确性部分回答“为什么结果可信”，成本部分回答“为什么这个实现会快或慢”。如果一个模型只能解释速度，却解释不了失败恢复，它不适合系统教材；如果只能解释语义，却完全不关心 cache、调度、I/O 或网络成本，也不适合第二册。

实验应围绕“只改变一个变量”设计。针对重试，最小实验可以保留朴素版本，再实现一个改进版本，输入规模从小到大，线程数或并发度从一到目标机器上限，错误注入从无到有。每一组实验都应记录 reference 结果、核心指标、环境和命令。这样读者看到差异时，可以判断差异来自机制本身，而不是来自初始化、缓存冷热、调度迁移或文件系统状态。

观察手段是：注入响应丢失和服务端慢处理。观察不是为了堆工具名，而是为了回答一个具体假设。若假设是共享写导致扩展性下降，就应看线程数曲线、写入布局和 cache 相关指标；若假设是提交边界错误，就应做崩溃点矩阵；若假设是队列背压，就应看水位、等待和上游速率。工具输出必须回到假设，不要把截图或命令输出当成结论。

最后要讲边界：不可幂等操作不能盲目重试。高质量教材必须把反例放在正文里。读者需要知道什么时候该用这个机制，什么时候它会制造新问题，什么时候应该退回更简单方案。比如一个单线程工具没有并发共享，强行引入复杂运行时只会增加维护成本；一个没有恢复要求的临时分析脚本，也不必承担完整 manifest 协议。

把这一点落实到代码审查，可以要求每个涉及重试的改动都回答五个问题：朴素版本是什么，新增机制保护了哪个不变量，新增成本在哪里，如何回退，哪些测试证明边界。这五个问题比“用了某某技术”更能判断质量，也能防止团队在没有证据时追逐复杂实现。

\topic{深挖：幂等}

围绕“幂等”，先把它放回一个具体问题：重复执行同一个请求为什么不应改变最终结果。如果这个问题只在本机分布式模拟里出现，读者很容易把它当成局部技巧；但在driver 向 worker 发送任务并等待结果中，它会沿着输入、执行、同步、I/O 和提交一路传播。所以讲解时不能只写定义，而要让读者看到它怎样从一个小失败变成系统性约束。

朴素写法通常是“服务端每收到一次就执行一次”。这句话看似简单，却包含几个默认前提：状态只有一个拥有者，执行顺序可预测，资源足够近，失败不会穿过边界，观察结果能够代表真实成本。这些前提在教学小程序中可能成立，在driver 向 worker 发送任务并等待结果里却会被输入规模、线程交错、硬件拓扑或故障注入逐个打破。

失败信号是“响应丢失后重试会重复计数或重复写输出”。注意，这个失败不一定表现为崩溃。它也可能表现为吞吐不再增长、p99 抖动、重启后结果多了一份、某个分区长尾、CPU 使用率很高但有效工作很少。教材要把这些信号和机制绑定起来，否则读者只会背术语，遇到真实现象时仍然不知道从哪里下手。

更可靠的推导顺序是先写出不变量。对幂等来说，不变量不是一句口号，而是本章要保护的事实：哪些数据只能写一次，哪些状态可以重试，哪些结果可以被缓存，哪些成本必须摊销，哪些错误必须外显。只要不变量没有写清，后面的代码、实验和优化都没有稳定坐标。

然后才引入模型：幂等表记录 request id、状态和结果，使重复请求返回同一结果。这个模型应同时解释正确性和成本。正确性部分回答“为什么结果可信”，成本部分回答“为什么这个实现会快或慢”。如果一个模型只能解释速度，却解释不了失败恢复，它不适合系统教材；如果只能解释语义，却完全不关心 cache、调度、I/O 或网络成本，也不适合第二册。

实验应围绕“只改变一个变量”设计。针对幂等，最小实验可以保留朴素版本，再实现一个改进版本，输入规模从小到大，线程数或并发度从一到目标机器上限，错误注入从无到有。每一组实验都应记录 reference 结果、核心指标、环境和命令。这样读者看到差异时，可以判断差异来自机制本身，而不是来自初始化、缓存冷热、调度迁移或文件系统状态。

观察手段是：测试执行成功但响应丢失。观察不是为了堆工具名，而是为了回答一个具体假设。若假设是共享写导致扩展性下降，就应看线程数曲线、写入布局和 cache 相关指标；若假设是提交边界错误，就应做崩溃点矩阵；若假设是队列背压，就应看水位、等待和上游速率。工具输出必须回到假设，不要把截图或命令输出当成结论。

最后要讲边界：幂等记录本身也要和业务提交同边界。高质量教材必须把反例放在正文里。读者需要知道什么时候该用这个机制，什么时候它会制造新问题，什么时候应该退回更简单方案。比如一个单线程工具没有并发共享，强行引入复杂运行时只会增加维护成本；一个没有恢复要求的临时分析脚本，也不必承担完整 manifest 协议。

把这一点落实到代码审查，可以要求每个涉及幂等的改动都回答五个问题：朴素版本是什么，新增机制保护了哪个不变量，新增成本在哪里，如何回退，哪些测试证明边界。这五个问题比“用了某某技术”更能判断质量，也能防止团队在没有证据时追逐复杂实现。

\topic{深挖：部分失败}

围绕“部分失败”，先把它放回一个具体问题：客户端活着、服务端活着但网络断开时系统处于什么状态。如果这个问题只在真实线上服务里出现，读者很容易把它当成局部技巧；但在driver 向 worker 发送任务并等待结果中，它会沿着输入、执行、同步、I/O 和提交一路传播。所以讲解时不能只写定义，而要让读者看到它怎样从一个小失败变成系统性约束。

朴素写法通常是“把失败视为进程崩溃或成功两种”。这句话看似简单，却包含几个默认前提：状态只有一个拥有者，执行顺序可预测，资源足够近，失败不会穿过边界，观察结果能够代表真实成本。这些前提在教学小程序中可能成立，在driver 向 worker 发送任务并等待结果里却会被输入规模、线程交错、硬件拓扑或故障注入逐个打破。

失败信号是“真实系统存在消息丢失、延迟、分区和旧 leader”。注意，这个失败不一定表现为崩溃。它也可能表现为吞吐不再增长、p99 抖动、重启后结果多了一份、某个分区长尾、CPU 使用率很高但有效工作很少。教材要把这些信号和机制绑定起来，否则读者只会背术语，遇到真实现象时仍然不知道从哪里下手。

更可靠的推导顺序是先写出不变量。对部分失败来说，不变量不是一句口号，而是本章要保护的事实：哪些数据只能写一次，哪些状态可以重试，哪些结果可以被缓存，哪些成本必须摊销，哪些错误必须外显。只要不变量没有写清，后面的代码、实验和优化都没有稳定坐标。

然后才引入模型：部分失败意味着一方无法可靠知道另一方状态。这个模型应同时解释正确性和成本。正确性部分回答“为什么结果可信”，成本部分回答“为什么这个实现会快或慢”。如果一个模型只能解释速度，却解释不了失败恢复，它不适合系统教材；如果只能解释语义，却完全不关心 cache、调度、I/O 或网络成本，也不适合第二册。

实验应围绕“只改变一个变量”设计。针对部分失败，最小实验可以保留朴素版本，再实现一个改进版本，输入规模从小到大，线程数或并发度从一到目标机器上限，错误注入从无到有。每一组实验都应记录 reference 结果、核心指标、环境和命令。这样读者看到差异时，可以判断差异来自机制本身，而不是来自初始化、缓存冷热、调度迁移或文件系统状态。

观察手段是：用消息丢弃、重复、延迟和乱序模拟。观察不是为了堆工具名，而是为了回答一个具体假设。若假设是共享写导致扩展性下降，就应看线程数曲线、写入布局和 cache 相关指标；若假设是提交边界错误，就应做崩溃点矩阵；若假设是队列背压，就应看水位、等待和上游速率。工具输出必须回到假设，不要把截图或命令输出当成结论。

最后要讲边界：超时只是未知，不是远端已停止。高质量教材必须把反例放在正文里。读者需要知道什么时候该用这个机制，什么时候它会制造新问题，什么时候应该退回更简单方案。比如一个单线程工具没有并发共享，强行引入复杂运行时只会增加维护成本；一个没有恢复要求的临时分析脚本，也不必承担完整 manifest 协议。

把这一点落实到代码审查，可以要求每个涉及部分失败的改动都回答五个问题：朴素版本是什么，新增机制保护了哪个不变量，新增成本在哪里，如何回退，哪些测试证明边界。这五个问题比“用了某某技术”更能判断质量，也能防止团队在没有证据时追逐复杂实现。

\topic{深挖：流控}

围绕“流控”，先把它放回一个具体问题：worker 返回 busy 时 driver 应该怎么做。如果这个问题只在单线程小输入里出现，读者很容易把它当成局部技巧；但在driver 向 worker 发送任务并等待结果中，它会沿着输入、执行、同步、I/O 和提交一路传播。所以讲解时不能只写定义，而要让读者看到它怎样从一个小失败变成系统性约束。

朴素写法通常是“继续提交任务直到队列满”。这句话看似简单，却包含几个默认前提：状态只有一个拥有者，执行顺序可预测，资源足够近，失败不会穿过边界，观察结果能够代表真实成本。这些前提在教学小程序中可能成立，在driver 向 worker 发送任务并等待结果里却会被输入规模、线程交错、硬件拓扑或故障注入逐个打破。

失败信号是“重试和排队放大故障”。注意，这个失败不一定表现为崩溃。它也可能表现为吞吐不再增长、p99 抖动、重启后结果多了一份、某个分区长尾、CPU 使用率很高但有效工作很少。教材要把这些信号和机制绑定起来，否则读者只会背术语，遇到真实现象时仍然不知道从哪里下手。

更可靠的推导顺序是先写出不变量。对流控来说，不变量不是一句口号，而是本章要保护的事实：哪些数据只能写一次，哪些状态可以重试，哪些结果可以被缓存，哪些成本必须摊销，哪些错误必须外显。只要不变量没有写清，后面的代码、实验和优化都没有稳定坐标。

然后才引入模型：流控把接收方容量反馈给发送方。这个模型应同时解释正确性和成本。正确性部分回答“为什么结果可信”，成本部分回答“为什么这个实现会快或慢”。如果一个模型只能解释速度，却解释不了失败恢复，它不适合系统教材；如果只能解释语义，却完全不关心 cache、调度、I/O 或网络成本，也不适合第二册。

实验应围绕“只改变一个变量”设计。针对流控，最小实验可以保留朴素版本，再实现一个改进版本，输入规模从小到大，线程数或并发度从一到目标机器上限，错误注入从无到有。每一组实验都应记录 reference 结果、核心指标、环境和命令。这样读者看到差异时，可以判断差异来自机制本身，而不是来自初始化、缓存冷热、调度迁移或文件系统状态。

观察手段是：记录 in flight、busy 响应和退避时间。观察不是为了堆工具名，而是为了回答一个具体假设。若假设是共享写导致扩展性下降，就应看线程数曲线、写入布局和 cache 相关指标；若假设是提交边界错误，就应做崩溃点矩阵；若假设是队列背压，就应看水位、等待和上游速率。工具输出必须回到假设，不要把截图或命令输出当成结论。

最后要讲边界：流控不是错误，它是保护系统的协议。高质量教材必须把反例放在正文里。读者需要知道什么时候该用这个机制，什么时候它会制造新问题，什么时候应该退回更简单方案。比如一个单线程工具没有并发共享，强行引入复杂运行时只会增加维护成本；一个没有恢复要求的临时分析脚本，也不必承担完整 manifest 协议。

把这一点落实到代码审查，可以要求每个涉及流控的改动都回答五个问题：朴素版本是什么，新增机制保护了哪个不变量，新增成本在哪里，如何回退，哪些测试证明边界。这五个问题比“用了某某技术”更能判断质量，也能防止团队在没有证据时追逐复杂实现。

\topic{深挖：Trace Id}

围绕“Trace Id”，先把它放回一个具体问题：一次作业跨多个 worker 后如何定位慢在哪里。如果这个问题只在单机多线程里出现，读者很容易把它当成局部技巧；但在driver 向 worker 发送任务并等待结果中，它会沿着输入、执行、同步、I/O 和提交一路传播。所以讲解时不能只写定义，而要让读者看到它怎样从一个小失败变成系统性约束。

朴素写法通常是“每个进程各写自己的日志”。这句话看似简单，却包含几个默认前提：状态只有一个拥有者，执行顺序可预测，资源足够近，失败不会穿过边界，观察结果能够代表真实成本。这些前提在教学小程序中可能成立，在driver 向 worker 发送任务并等待结果里却会被输入规模、线程交错、硬件拓扑或故障注入逐个打破。

失败信号是“无法串联同一请求的路径”。注意，这个失败不一定表现为崩溃。它也可能表现为吞吐不再增长、p99 抖动、重启后结果多了一份、某个分区长尾、CPU 使用率很高但有效工作很少。教材要把这些信号和机制绑定起来，否则读者只会背术语，遇到真实现象时仍然不知道从哪里下手。

更可靠的推导顺序是先写出不变量。对Trace Id来说，不变量不是一句口号，而是本章要保护的事实：哪些数据只能写一次，哪些状态可以重试，哪些结果可以被缓存，哪些成本必须摊销，哪些错误必须外显。只要不变量没有写清，后面的代码、实验和优化都没有稳定坐标。

然后才引入模型：trace id 把跨组件事件连成一条时间线。这个模型应同时解释正确性和成本。正确性部分回答“为什么结果可信”，成本部分回答“为什么这个实现会快或慢”。如果一个模型只能解释速度，却解释不了失败恢复，它不适合系统教材；如果只能解释语义，却完全不关心 cache、调度、I/O 或网络成本，也不适合第二册。

实验应围绕“只改变一个变量”设计。针对Trace Id，最小实验可以保留朴素版本，再实现一个改进版本，输入规模从小到大，线程数或并发度从一到目标机器上限，错误注入从无到有。每一组实验都应记录 reference 结果、核心指标、环境和命令。这样读者看到差异时，可以判断差异来自机制本身，而不是来自初始化、缓存冷热、调度迁移或文件系统状态。

观察手段是：记录 driver send、worker receive、execute、reply。观察不是为了堆工具名，而是为了回答一个具体假设。若假设是共享写导致扩展性下降，就应看线程数曲线、写入布局和 cache 相关指标；若假设是提交边界错误，就应做崩溃点矩阵；若假设是队列背压，就应看水位、等待和上游速率。工具输出必须回到假设，不要把截图或命令输出当成结论。

最后要讲边界：trace 也要采样和限流，不能让观测拖垮系统。高质量教材必须把反例放在正文里。读者需要知道什么时候该用这个机制，什么时候它会制造新问题，什么时候应该退回更简单方案。比如一个单线程工具没有并发共享，强行引入复杂运行时只会增加维护成本；一个没有恢复要求的临时分析脚本，也不必承担完整 manifest 协议。

把这一点落实到代码审查，可以要求每个涉及Trace Id的改动都回答五个问题：朴素版本是什么，新增机制保护了哪个不变量，新增成本在哪里，如何回退，哪些测试证明边界。这五个问题比“用了某某技术”更能判断质量，也能防止团队在没有证据时追逐复杂实现。

\topic{深挖：本机模拟}

围绕“本机模拟”，先把它放回一个具体问题：没有集群时怎样学习分布式失败。如果这个问题只在NUMA 或异构核心里出现，读者很容易把它当成局部技巧；但在driver 向 worker 发送任务并等待结果中，它会沿着输入、执行、同步、I/O 和提交一路传播。所以讲解时不能只写定义，而要让读者看到它怎样从一个小失败变成系统性约束。

朴素写法通常是“只写正常路径函数调用”。这句话看似简单，却包含几个默认前提：状态只有一个拥有者，执行顺序可预测，资源足够近，失败不会穿过边界，观察结果能够代表真实成本。这些前提在教学小程序中可能成立，在driver 向 worker 发送任务并等待结果里却会被输入规模、线程交错、硬件拓扑或故障注入逐个打破。

失败信号是“无法暴露乱序、重复和超时”。注意，这个失败不一定表现为崩溃。它也可能表现为吞吐不再增长、p99 抖动、重启后结果多了一份、某个分区长尾、CPU 使用率很高但有效工作很少。教材要把这些信号和机制绑定起来，否则读者只会背术语，遇到真实现象时仍然不知道从哪里下手。

更可靠的推导顺序是先写出不变量。对本机模拟来说，不变量不是一句口号，而是本章要保护的事实：哪些数据只能写一次，哪些状态可以重试，哪些结果可以被缓存，哪些成本必须摊销，哪些错误必须外显。只要不变量没有写清，后面的代码、实验和优化都没有稳定坐标。

然后才引入模型：本机 message bus 可以注入延迟、丢失、重复和乱序。这个模型应同时解释正确性和成本。正确性部分回答“为什么结果可信”，成本部分回答“为什么这个实现会快或慢”。如果一个模型只能解释速度，却解释不了失败恢复，它不适合系统教材；如果只能解释语义，却完全不关心 cache、调度、I/O 或网络成本，也不适合第二册。

实验应围绕“只改变一个变量”设计。针对本机模拟，最小实验可以保留朴素版本，再实现一个改进版本，输入规模从小到大，线程数或并发度从一到目标机器上限，错误注入从无到有。每一组实验都应记录 reference 结果、核心指标、环境和命令。这样读者看到差异时，可以判断差异来自机制本身，而不是来自初始化、缓存冷热、调度迁移或文件系统状态。

观察手段是：用固定随机种子复现实验。观察不是为了堆工具名，而是为了回答一个具体假设。若假设是共享写导致扩展性下降，就应看线程数曲线、写入布局和 cache 相关指标；若假设是提交边界错误，就应做崩溃点矩阵；若假设是队列背压，就应看水位、等待和上游速率。工具输出必须回到假设，不要把截图或命令输出当成结论。

最后要讲边界：模拟不能证明真实网络性能，但能训练协议语义。高质量教材必须把反例放在正文里。读者需要知道什么时候该用这个机制，什么时候它会制造新问题，什么时候应该退回更简单方案。比如一个单线程工具没有并发共享，强行引入复杂运行时只会增加维护成本；一个没有恢复要求的临时分析脚本，也不必承担完整 manifest 协议。

把这一点落实到代码审查，可以要求每个涉及本机模拟的改动都回答五个问题：朴素版本是什么，新增机制保护了哪个不变量，新增成本在哪里，如何回退，哪些测试证明边界。这五个问题比“用了某某技术”更能判断质量，也能防止团队在没有证据时追逐复杂实现。

\topic{案例深化：从朴素版到工程版}

案例深化“响应丢失”要按三次迭代写。第一次只写 reference：worker 已经完成任务并提交结果，但回复在路上丢失。reference 的代码可以慢，但必须把输入、输出、错误和边界写清。第二次写朴素工程版本：driver 超时后派发新 attempt。这一步不能省，因为读者需要看到直觉方案为什么诱人，也需要有一个失败对象。

第三次才写改进版本：worker 或 reduce 端用 request id 去重。改进说明要避免“因为更高级所以更快”这种表达，而要讲清它改变了哪条路径：减少了数据移动、减少了共享写、减少了系统调用、减少了重试放大，还是把不可恢复状态变成可恢复状态。

验收方式是：结果只提交一次，重复响应返回同一摘要。验收报告至少包含一张对照表，一列是朴素版本，一列是改进版本，一列是反例。反例很重要，因为它能告诉读者改进不是什么时候都正确。如果一个案例没有反例，往往说明分析还停在宣传层面。
案例深化“重试风暴”要按三次迭代写。第一次只写 reference：服务端变慢，客户端同时超时并重试。reference 的代码可以慢，但必须把输入、输出、错误和边界写清。第二次写朴素工程版本：所有客户端立即重试。这一步不能省，因为读者需要看到直觉方案为什么诱人，也需要有一个失败对象。

第三次才写改进版本：指数退避、抖动和全局 in flight 限制。改进说明要避免“因为更高级所以更快”这种表达，而要讲清它改变了哪条路径：减少了数据移动、减少了共享写、减少了系统调用、减少了重试放大，还是把不可恢复状态变成可恢复状态。

验收方式是：请求量不会因重试超过预算。验收报告至少包含一张对照表，一列是朴素版本，一列是改进版本，一列是反例。反例很重要，因为它能告诉读者改进不是什么时候都正确。如果一个案例没有反例，往往说明分析还停在宣传层面。
案例深化“乱序响应”要按三次迭代写。第一次只写 reference：attempt 2 的响应先到，attempt 1 的响应后到。reference 的代码可以慢，但必须把输入、输出、错误和边界写清。第二次写朴素工程版本：谁后到就覆盖状态。这一步不能省，因为读者需要看到直觉方案为什么诱人，也需要有一个失败对象。

第三次才写改进版本：用 attempt 和 epoch 做状态检查。改进说明要避免“因为更高级所以更快”这种表达，而要讲清它改变了哪条路径：减少了数据移动、减少了共享写、减少了系统调用、减少了重试放大，还是把不可恢复状态变成可恢复状态。

验收方式是：旧 attempt 不会覆盖已提交结果。验收报告至少包含一张对照表，一列是朴素版本，一列是改进版本，一列是反例。反例很重要，因为它能告诉读者改进不是什么时候都正确。如果一个案例没有反例，往往说明分析还停在宣传层面。

\topic{实验矩阵：让结论可复现}

围绕消息身份，实验矩阵至少包含正常输入、边界输入、压力输入和错误输入四类。正常输入验证功能，边界输入验证不变量，压力输入验证成本模型，错误输入验证恢复和报告。其中压力输入不只是“大”，还要能触发本章关心的形状：随机访问、热点 key、短任务、慢 I/O、重复消息或跨节点访问。

围绕Deadline，实验矩阵至少包含正常输入、边界输入、压力输入和错误输入四类。正常输入验证功能，边界输入验证不变量，压力输入验证成本模型，错误输入验证恢复和报告。其中压力输入不只是“大”，还要能触发本章关心的形状：随机访问、热点 key、短任务、慢 I/O、重复消息或跨节点访问。

围绕重试，实验矩阵至少包含正常输入、边界输入、压力输入和错误输入四类。正常输入验证功能，边界输入验证不变量，压力输入验证成本模型，错误输入验证恢复和报告。其中压力输入不只是“大”，还要能触发本章关心的形状：随机访问、热点 key、短任务、慢 I/O、重复消息或跨节点访问。

围绕幂等，实验矩阵至少包含正常输入、边界输入、压力输入和错误输入四类。正常输入验证功能，边界输入验证不变量，压力输入验证成本模型，错误输入验证恢复和报告。其中压力输入不只是“大”，还要能触发本章关心的形状：随机访问、热点 key、短任务、慢 I/O、重复消息或跨节点访问。

围绕部分失败，实验矩阵至少包含正常输入、边界输入、压力输入和错误输入四类。正常输入验证功能，边界输入验证不变量，压力输入验证成本模型，错误输入验证恢复和报告。其中压力输入不只是“大”，还要能触发本章关心的形状：随机访问、热点 key、短任务、慢 I/O、重复消息或跨节点访问。

围绕流控，实验矩阵至少包含正常输入、边界输入、压力输入和错误输入四类。正常输入验证功能，边界输入验证不变量，压力输入验证成本模型，错误输入验证恢复和报告。其中压力输入不只是“大”，还要能触发本章关心的形状：随机访问、热点 key、短任务、慢 I/O、重复消息或跨节点访问。

围绕Trace Id，实验矩阵至少包含正常输入、边界输入、压力输入和错误输入四类。正常输入验证功能，边界输入验证不变量，压力输入验证成本模型，错误输入验证恢复和报告。其中压力输入不只是“大”，还要能触发本章关心的形状：随机访问、热点 key、短任务、慢 I/O、重复消息或跨节点访问。

围绕本机模拟，实验矩阵至少包含正常输入、边界输入、压力输入和错误输入四类。正常输入验证功能，边界输入验证不变量，压力输入验证成本模型，错误输入验证恢复和报告。其中压力输入不只是“大”，还要能触发本章关心的形状：随机访问、热点 key、短任务、慢 I/O、重复消息或跨节点访问。

\topic{Linux 源码阅读深化}

本章的 Linux 阅读入口仍然保持窄范围：\filepath{net/core/dev.c}、\filepath{net/ipv4/tcp.c}、\filepath{kernel/time/timer.c}。阅读时不要追求覆盖率，而要追求因果链。从一个用户态现象出发，找到内核中的状态对象，再看状态怎样被修改、等待怎样被挂起、唤醒怎样发生、错误怎样返回。

源码阅读可以按四栏笔记整理。第一栏写用户态操作，例如线程等待、缺页、写文件、发送消息或计时。第二栏写内核对象，例如 task、VMA、page、inode、wait queue、bio、socket buffer。第三栏写关键状态变化，例如 runnable 到 sleeping、clean page 到 dirty page、pending task 到 committed task。第四栏写可观察指标或实验，例如 context switch、page fault、writeback、queue depth、retry count。

不要把 Linux 源码当成术语来源。源码阅读的目标是训练映射能力：把书中的抽象边界映射到真实系统里的结构和状态。读者不需要一次读懂整个内核，但要能解释一个最小现象为什么发生、哪些路径参与、哪些结论受当前内核版本和硬件限制。

\topic{项目验收：把本章能力写进仓库}

贯穿项目的验收不能只说“功能跑通”。本章至少要留下一个可重复实验、一个失败注入、一个指标输出和一个设计说明。实验说明机制是否真的被触发；失败注入说明边界是否真的被处理；指标输出说明结论是否有证据；设计说明让后续章节能复用这个模块。

本章项目检查点可以具体化为：1. 为每条消息加入 request id 和 attempt
2. 实现可注入故障的 message bus
3. 加入 deadline 和重试预算
4. 实现幂等结果表
5. 输出跨组件 trace。这些检查点要进入仓库，而不是停留在阅读笔记。如果某个检查点因为当前设备权限、硬件或系统 API 不支持而无法完成，也要写清降级方案和未验证风险。

项目验收还应要求读者写一段复盘：本章新增机制让系统多了什么能力，也让系统多了什么复杂度。比如加入背压后内存更稳定，但生产者会等待；加入 route epoch 后旧请求能被拒绝，但所有消息都必须携带版本；加入 SIMD 后热内核更快，但尾部和数值语义要额外测试。这种复盘能把知识从“会用”推进到“会判断”。



\topic{最终收束：本章应形成的判断力}

读完本章，读者不应只记住消息身份、本机模拟这些词，而应能围绕driver 向 worker 发送任务并等待结果做一次完整判断。完整判断从问题开始：现象是什么，朴素方案是什么，失败信号是什么，保护的不变量是什么，成本从哪里移动到哪里。若无法回答这些问题，就说明还停留在概念层，没有真正进入系统层。

本章最重要的反例是：远端执行的完成、失败和重复都无法靠本地调用栈表达。遇到类似反例时，第一步不是换技术，而是缩小问题。先固定输入，固定环境，保留 reference，再一次只改变一个变量。如果改动同时改变布局、线程数、批大小、同步方式和提交策略，即使结果变快，也无法知道原因。高质量工程实验必须克制变量。

以案例“响应丢失”为例，最终报告应包含三段。第一段写worker 已经完成任务并提交结果，但回复在路上丢失，说明读者要解决的不是抽象题，而是一个有输入、有状态、有失败方式的任务。第二段写朴素版本：driver 超时后派发新 attempt，并诚实说明它为什么吸引人。第三段写改进版本：worker 或 reduce 端用 request id 去重，再用结果只提交一次，重复响应返回同一摘要验收。报告中若没有朴素版本，读者会误以为最终设计是凭空出现；若没有反例，读者会误以为最终设计永远正确。

本章还应留下一个审查表。正确性方面，检查输入合同、状态转移、所有权、重复执行、关闭和恢复。性能方面，检查数据移动、共享写、排队、系统调用、批大小、缓存冷热和拓扑。可观测性方面，检查是否有阶段耗时、队列水位、错误分类、任务身份和原始样本。每一项都要能落到代码、测试或报告，不要只停在口头承诺。

贯穿项目里，本章至少要完成这些检查点：为每条消息加入 request id 和 attempt、实现可注入故障的 message bus、加入 deadline 和重试预算、实现幂等结果表、输出跨组件 trace。完成后不要急着进入下一章，要先写一页复盘：本章新增能力解决了什么失败，新增了哪些复杂度，哪些结论只在当前机器上验证，哪些结论可以迁移到更大系统。这种复盘会让整本书的知识保持连续，不会变成每章讲完就散掉的材料。

最后，用一句话收束本章：系统能力来自清楚的边界、可验证的不变量和可复现的证据。无论本章讨论的是硬件执行、内存层级、同步、运行时、I/O 还是分布式协议，只要缺少边界、不变量和证据，复杂实现就会变成风险来源。反过来，只要这三件事稳住，读者就能把一个朴素程序逐步推进成可靠的计算系统。




\topic{过线补充：常见误区、验收题和迁移能力}

本章最容易出现的误区，是把消息身份、Deadline、重试、幂等当成互相独立的知识点。在driver 向 worker 发送任务并等待结果中，它们其实围绕同一个失败展开：远端执行的完成、失败和重复都无法靠本地调用栈表达。如果读者只能分别解释这些概念，却不能说明它们在同一条数据流里怎样互相影响，就还没有真正掌握本章。例如一个优化减少了计算时间，却增加了队列等待；一个同步方案修复了正确性，却制造了共享写热点；一个提交协议提高了可靠性，却增加了持久化延迟。这些都不是矛盾，而是系统设计中的成本迁移。

本章的验收题可以这样设计：先给出一个最小版本的driver 向 worker 发送任务并等待结果，要求读者写出 reference；再引入一个压力条件，要求读者指出朴素方案会在哪里失败；最后要求读者选择本章机制中的一项或两项做改进。答案不能只给最终代码，还要给不变量、实验矩阵和反例。如果某个答案没有说明失败后如何恢复、如何观察、如何判断优化是否值得，就不能算完整答案。

围绕案例响应丢失、重试风暴、乱序响应，报告还应补一个迁移问题：同样方法迁移到更大输入、更多线程、不同硬件或本机分布式模拟时，哪些结论保持，哪些结论要重新测。保持的通常是语义边界和不变量；需要重测的通常是吞吐、延迟、批大小、缓存效果、锁竞争和 I/O 行为。这能训练读者区分“原理可以迁移”和“数字不能照搬”。

最后再强调一次本书的写法要求：先从问题出发，再推导机制，再落到实验。不要把实验放成装饰，也不要把 Linux 源码阅读放成资料链接。实验必须检验一个假设，源码阅读必须解释一个用户态现象，项目代码必须保护一个明确边界。读者能按这个顺序复述本章，才说明本章内容真正连贯起来。




\topic{事故复盘式走查}

为了把本章内容真正串起来，可以把driver 向 worker 发送任务并等待结果写成一次小型事故复盘。事故现象不是一句“系统慢了”或“结果错了”，而要具体到可观察事实：哪一批输入、哪个阶段、哪一个指标、哪一种失败注入触发了问题。如果现象描述不具体，后面的消息身份、Deadline和重试就只能变成猜测。

复盘第一步是还原朴素设计。写清当时为什么选择它，它解决了什么简单问题，又默认了哪些条件。很多坏设计一开始并不坏，只是从小输入、小并发、无失败的条件迁移到driver 向 worker 发送任务并等待结果时，没有同步升级边界。这也是本册反复强调从朴素方案出发的原因：只有理解朴素方案为什么成立，才能准确说明它为什么在新条件下失败。

复盘第二步是列出候选假设，但每个假设都必须有可证伪实验。如果怀疑消息身份相关路径，就构造只改变这一因素的对照；如果怀疑Deadline，就记录它对应的状态或指标；如果怀疑重试，就找出它在代码里的所有入口和退出点。不能把所有怀疑同时改掉。一次修改多个因素会让复盘变成经验叙事，而不是工程证据。

复盘第三步是修复。修复不等于把代码改到通过测试，而是把不变量写回系统。本章的不变量来自“远端执行的完成、失败和重复都无法靠本地调用栈表达”这个失败主题。因此修复说明要写清：新增字段保护什么，新增队列容量限制什么，新增版本号拒绝什么，新增 reference 比较什么，新增指标暴露什么。如果修复无法对应到不变量，就很可能只是局部补丁。

复盘第四步是验收。以“乱序响应”为例，验收不应只跑正常输入，还要跑边界输入、压力输入和错误输入。正常输入证明功能还在，边界输入证明不变量没有被破坏，压力输入证明成本模型仍然成立，错误输入证明失败不会越过提交边界。验收报告要保留原始命令和数据，否则下一次回归无法判断问题是否真正修复。

最后要写“未解决问题”。例如本机模拟相关边界可能还没有在当前设备上完全验证，某些 Linux 计数器可能因为权限不可用，某些 NUMA 结论可能因为硬件不支持只能停留在设计层。把未验证风险写出来不是降低质量，而是防止教材把假设写成事实。高质量书籍要敢于说明证据边界，这比把每个结论都写成绝对经验更可靠。

读者完成本章后，可以用这份复盘模板检查自己的学习结果：能否描述事故，能否还原朴素设计，能否提出可证伪假设，能否把修复绑定到不变量，能否设计验收矩阵，能否写出未验证风险。如果六项都能完成，本章知识就已经从概念记忆转成工程判断。如果只能完成其中一两项，就应回到本章前面的案例重新走一遍，而不是急着进入下一章。



\section{本章落到贯穿项目}

Compute Systems Engine 的分布式模拟从三个对象开始：driver、worker、message bus。driver 生成任务和 request id；worker 读取本地分片并执行 map/reduce；message bus 用有界队列模拟网络，支持随机延迟、超时、重复和丢弃。每条消息携带 request id、attempt、source、target partition 和 payload size。

第一版只做日志计数：worker 本地统计 key，再把 partial 发给 reduce partition。第二版加入重试：message bus 随机丢响应，driver 重试同一 request id，reduce 端必须去重。第三版加入背压：reduce 队列满时返回 busy 或延迟，map 端降低发送速率。第四版加入检查点：每个 partition 记录已提交 request id 和结果。

\section{本章习题}

\begin{exercisebox}
习题一：画出一次 RPC 的完整生命周期，从客户端提交到服务端响应返回。标出每个排队点、失败点和需要记录的指标。

习题二：设计一个幂等请求表。说明 request id、业务状态、响应结果如何在同一个提交边界内更新。给出崩溃发生在提交前和提交后的恢复行为。

习题三：构造一个重试风暴。假设服务端延迟升高，客户端超时后立即重试。计算请求量如何放大，并设计重试预算、退避和抖动策略。

习题四：用本机多线程模拟消息乱序和重复。要求每条消息携带 request id 和 attempt，服务端能识别重复请求并返回同一结果。

习题五：比较数据本地调度和远程调度。给出本地排队长、远程网络快以及本地网络快、远程排队短两种场景，说明调度器应如何取舍。
\end{exercisebox}

\begin{labbox}
实验：在本机用多个进程模拟客户端和服务器。人为加入随机延迟和超时，观察重试如何造成重复请求。设计一个幂等请求 ID 机制。
\end{labbox}
